% =====================================================================
%  RELATÓRIO DE PROJETO FINAL DE LICENCIATURA
%  Sistema de Geração Automática de Histórias Familiares com IA
%  Universidade da Beira Interior — Departamento de Informática
%
%  Documento mestre. Compila em Overleaf com pdfLaTeX + BibTeX.
%  Cada capítulo está num ficheiro próprio (ver \input mais abaixo).
%
%  COMO COMPILAR NO OVERLEAF:
%    1. Carregar TODA a pasta "Projetos" (incluindo a subpasta images/).
%    2. Definir main.tex como documento principal.
%    3. Compilador: pdfLaTeX. Premir "Recompile" duas vezes para que
%       o índice, as referências cruzadas e a bibliografia estabilizem.
%
%  IMAGENS: ver o ficheiro IMAGENS_NECESSARIAS.md. Onde falta uma
%  figura, aparece uma moldura "[Imagem a inserir]". Para inserir a
%  imagem real, troca o comando \figph{...} por \includegraphics.
% =====================================================================
\documentclass[12pt,a4paper,oneside]{report}

% ---- Codificação e língua ------------------------------------------
\usepackage[utf8]{inputenc}
\usepackage[T1]{fontenc}
\usepackage[portuguese]{babel}
\usepackage{lmodern}
\usepackage{microtype}

% ---- Layout e espaçamento ------------------------------------------
\usepackage[a4paper,top=2.5cm,bottom=2.5cm,left=3cm,right=2.5cm]{geometry}
\usepackage{setspace}
\onehalfspacing
\usepackage{emptypage}
\usepackage{indentfirst}

% ---- Gráficos, tabelas, listas -------------------------------------
\usepackage{graphicx}
\graphicspath{{images/}}
\usepackage{float}
\usepackage{booktabs}
\usepackage{tabularx}
\usepackage{longtable}
\usepackage{array}
\usepackage{multirow}
\usepackage{enumitem}
\usepackage{caption}
\usepackage{subcaption}
\captionsetup{font=small,labelfont=bf}

% ---- Cor -----------------------------------------------------------
\usepackage{xcolor}
\definecolor{ubiblue}{RGB}{0,45,98}
\definecolor{accent}{RGB}{150,90,20}
\definecolor{codebg}{RGB}{246,246,246}
\definecolor{codekw}{RGB}{160,50,40}
\definecolor{codecm}{RGB}{90,110,90}
\definecolor{codestr}{RGB}{30,110,70}

% ---- Código (listings) ---------------------------------------------
\usepackage{listings}
\renewcommand{\lstlistingname}{Excerto}
\renewcommand{\lstlistlistingname}{Lista de Excertos de Código}
\lstdefinestyle{projeto}{
  backgroundcolor=\color{codebg},
  basicstyle=\ttfamily\footnotesize,
  keywordstyle=\color{codekw}\bfseries,
  commentstyle=\color{codecm}\itshape,
  stringstyle=\color{codestr},
  breaklines=true,
  showstringspaces=false,
  frame=single,
  rulecolor=\color{gray!35},
  numbers=left,
  numberstyle=\tiny\color{gray},
  tabsize=2,
  captionpos=b,
  columns=fullflexible,
  xleftmargin=16pt,
  framexleftmargin=14pt,
  % Permite acentos portugueses (UTF-8) nos comentários das listagens.
  extendedchars=true,
  literate=%
    {á}{{\'a}}1 {à}{{\`a}}1 {â}{{\^a}}1 {ã}{{\~a}}1
    {é}{{\'e}}1 {ê}{{\^e}}1 {í}{{\'i}}1 {ó}{{\'o}}1
    {ô}{{\^o}}1 {õ}{{\~o}}1 {ú}{{\'u}}1 {ç}{{\c c}}1
    {Á}{{\'A}}1 {É}{{\'E}}1 {Í}{{\'I}}1 {Ó}{{\'O}}1
    {Ú}{{\'U}}1 {Ã}{{\~A}}1 {Õ}{{\~O}}1 {Ç}{{\c C}}1
    {Ê}{{\^E}}1 {Â}{{\^A}}1 {Ô}{{\^O}}1 {À}{{\`A}}1,
}
\lstset{style=projeto}

% ---- TikZ (diagramas) ----------------------------------------------
\usepackage{tikz}
\usetikzlibrary{shapes.geometric,arrows.meta,positioning,fit,backgrounds,calc}

% ---- Cabeçalhos/rodapés --------------------------------------------
\usepackage{fancyhdr}
\pagestyle{fancy}
\fancyhf{}
\fancyhead[L]{\small\nouppercase{\leftmark}}
\fancyhead[R]{\small\thepage}
\renewcommand{\headrulewidth}{0.4pt}
\fancypagestyle{plain}{\fancyhf{}\fancyfoot[C]{\small\thepage}\renewcommand{\headrulewidth}{0pt}}

% ---- Hiperligações (carregar quase no fim) -------------------------
\usepackage[hidelinks]{hyperref}
\hypersetup{
  colorlinks=true,
  linkcolor=ubiblue,
  citecolor=accent,
  urlcolor=accent,
  pdftitle={Geração Automática de Histórias Familiares com IA},
  pdfauthor={Diogo Dinis},
}
\usepackage{cleveref}

% ---- Comandos auxiliares -------------------------------------------
% Placeholder de figura: desenha uma moldura tracejada com a legenda
% do que deve ser inserido. Compila sem precisar do ficheiro de imagem.
\newcommand{\figph}[2][9cm]{%
  \begin{center}
  \setlength{\fboxsep}{10pt}%
  \fcolorbox{gray!50}{gray!8}{\parbox[c][#1][c]{0.86\linewidth}{%
    \centering\itshape\color{gray!60!black}%
    [\,IMAGEM A INSERIR\,]\\[6pt]\normalfont\small #2}}%
  \end{center}}

\newcolumntype{Y}{>{\raggedright\arraybackslash}X}
\newcommand{\rf}[1]{\textbf{RF#1}}
\newcommand{\rnf}[1]{\textbf{RNF#1}}

% =====================================================================
\begin{document}

% ----------------------------- CAPA ---------------------------------
% =====================================================================
%  CAPA / FOLHA DE ROSTO
%  Estilo Universidade da Beira Interior — Faculdade de Engenharia.
%
%  LOGÓTIPO: coloca o ficheiro do logótipo oficial em
%      images/ubi-fe-di.png
%  Assim que o ficheiro existir, ele é usado automaticamente; enquanto
%  não existir, é mostrado um placeholder (o documento compila na mesma).
% =====================================================================
\begin{titlepage}
\centering
\thispagestyle{empty}

\vspace*{0.6cm}
{\Huge\bfseries Universidade da Beira Interior}\\[0.35cm]
{\Huge Departamento de Informática}\\[1.1cm]

% --- Logótipo institucional (FE-UBI / Departamento de Informática) ---
\IfFileExists{ubi-fe-di.png}{%
  \includegraphics[width=191pt]{ubi-fe-di.png}%
}{%
  \figph[5cm]{Logótipo oficial da Faculdade de Engenharia da UBI
  (coloca o ficheiro \texttt{images/ubi-fe-di.png}).}%
}

\vspace{1.2cm}

{\color{ubiblue}\rule{\linewidth}{1pt}}\\[0.35cm]
{\huge\bfseries Sistema de Geração Automática de Histórias
Familiares com Inteligência Artificial Generativa}\\[0.35cm]
{\color{ubiblue}\rule{\linewidth}{1pt}}\\[0.9cm]

{\Large\bfseries Relatório de Projeto Final de Licenciatura}\\[0.25cm]
{\large Licenciatura em Engenharia Informática}\\[1.1cm]

{\large\bfseries Elaborado por:}\\[0.25cm]
{\large Diogo Dinis, n.º 52374}\\[1.1cm]

{\large\bfseries Orientadores:}\\[0.25cm]
{\large Professor Vasco Lopes}\\
{\large Professor João Dias}

\vfill
{\large Covilhã, \today}
\vspace*{0.4cm}
\end{titlepage}


% ------------------------ ELEMENTOS PRÉ-TEXTUAIS --------------------
\pagenumbering{roman}
\setcounter{page}{2}

% =====================================================================
%  RESUMO + ABSTRACT
% =====================================================================
\chapter*{Resumo}
\addcontentsline{toc}{chapter}{Resumo}
\markboth{Resumo}{Resumo}

O arquivo familiar — fotografias, documentos e árvores genealógicas
acumulados ao longo de gerações — encerra um enorme valor afetivo e
histórico, mas permanece tipicamente disperso, desorganizado e, acima de
tudo, \emph{não narrado}. Este projeto propõe e implementa um sistema que
transforma automaticamente esse material num produto narrativo
multimédia, sem exigir que o utilizador escreva textos ou descreva o
conteúdo manualmente.

O sistema organiza-se num \emph{pipeline} de quatro módulos sequenciais:
\textbf{(M1)} ingestão multimodal, que extrai metadados \textsc{exif},
reconhece texto em documentos por \textsc{ocr} e descreve o conteúdo
visual das fotografias com um modelo de visão; \textbf{(M2)} organização
temporal, que resolve datas incompletas e constrói um grafo de
parentesco; \textbf{(M3)} geração narrativa, assente em modelos de
linguagem de grande escala (\textsc{llm}) orientados por
\emph{Retrieval-Augmented Generation} (\textsc{rag}) e por um conjunto de
\emph{templates} calibrados para português europeu; e \textbf{(M4)}
geração multimédia, que converte a narrativa num vídeo documental com
enquadramento cinematográfico, narração por síntese de voz e legendas
sincronizadas. A narrativa é segmentada em \emph{cenas} ligadas a
fotografias concretas, o que permite ao M4 sincronizar cada imagem com a
sua narração.

A solução foi desenvolvida como uma aplicação Web completa, com um
\emph{backend} assíncrono em \textsc{FastAPI} e um \emph{frontend} em
\textsc{React}/\textsc{TypeScript}, e evoluiu de um protótipo
\emph{local-first} para uma arquitetura cliente-servidor multi-inquilino
implantada na nuvem (\textsc{Render}, \textsc{Vercel} e \textsc{Supabase}
para autenticação, base de dados e armazenamento), correndo de forma
idêntica em ambiente local e \emph{online}. Adota uma estratégia de
inteligência artificial \textbf{em cascata e a custo nulo} — texto via
\emph{Groq} (Llama~3.3~70B, na nuvem), com a inferência local via
\emph{Ollama} e o \emph{Google~Gemini} como alternativas, e visão via
\emph{Gemini} — e técnicas de engenharia de robustez (degradação graciosa
do \textsc{rag}, \emph{fallbacks} de modelo e rebaixamento síncrono) que
mantêm o sistema operacional mesmo em ambientes com dependências
reduzidas. Os resultados
demonstram a viabilidade de gerar, de forma autónoma e a custo
essencialmente nulo, narrativas factualmente ancoradas e vídeos
documentais coerentes a partir de um arquivo familiar heterogéneo.

\vspace{0.8cm}
\noindent\textbf{Palavras-chave}\\[0.2cm]
Inteligência Artificial Generativa; Modelos de Linguagem de Grande
Escala; Retrieval-Augmented Generation; Processamento Multimodal;
Genealogia Computacional; Geração Automática de Vídeo; Narrativa Digital;
FastAPI; React.

\chapter*{Abstract}
\addcontentsline{toc}{chapter}{Abstract}
\markboth{Abstract}{Abstract}

The family archive — photographs, documents and genealogical trees
accumulated across generations — holds enormous emotional and historical
value, yet it typically remains scattered, disorganised and, above all,
\emph{untold}. This project proposes and implements a system that
automatically turns such material into a multimedia narrative product,
without requiring the user to write any text or describe the content
manually.

The system is organised as a pipeline of four sequential modules:
\textbf{(M1)} multimodal ingestion, extracting \textsc{exif} metadata,
running \textsc{ocr} on documents and describing the visual content of
photographs with a vision model; \textbf{(M2)} temporal organisation,
resolving incomplete dates and building a kinship graph; \textbf{(M3)}
narrative generation, based on Large Language Models steered by
Retrieval-Augmented Generation and a set of European-Portuguese prompt
templates; and \textbf{(M4)} multimedia generation, turning the narrative
into a documentary video with cinematic framing, text-to-speech
narration and synchronised captions.

The solution was built as a complete web application, with an
asynchronous \textsc{FastAPI} backend and a \textsc{React}/\textsc{TypeScript}
frontend, evolving from a local-first prototype into a cloud-deployed,
multi-tenant client–server architecture (Render, Vercel and Supabase for
authentication, database and storage), running identically in local and
online environments. It adopts a \textbf{zero-cost, cascading AI strategy}
— text via Groq (cloud Llama~3.3~70B), with local Ollama inference and
Google~Gemini as fallbacks, and vision via Gemini — and robustness
engineering techniques (graceful RAG degradation, model fallbacks and
synchronous downgrade) that keep the system operational even in
dependency-constrained environments. The
results demonstrate the feasibility of autonomously generating, at
essentially zero cost, factually grounded narratives and coherent
documentary videos from a heterogeneous family archive.

\vspace{0.8cm}
\noindent\textbf{Keywords}\\[0.2cm]
Generative Artificial Intelligence; Large Language Models;
Retrieval-Augmented Generation; Multimodal Processing; Computational
Genealogy; Automatic Video Generation; Digital Storytelling; FastAPI;
React.

% =====================================================================
%  AGRADECIMENTOS (opcional — adapta livremente o texto)
% =====================================================================
\chapter*{Agradecimentos}
\addcontentsline{toc}{chapter}{Agradecimentos}
\markboth{Agradecimentos}{Agradecimentos}

Aos meus orientadores, \textbf{Professor Vasco Lopes} e \textbf{Professor
João Dias}, pela orientação, disponibilidade e visão crítica que deram
forma a este projeto ao longo do seu desenvolvimento.

À \textbf{Universidade da Beira Interior} e ao \textbf{Departamento de
Informática}, pelo ambiente académico e pelos recursos que tornaram este
trabalho possível.

À minha \textbf{família}, cujo arquivo de memórias inspirou o próprio
tema deste projeto, e cujo apoio incondicional acompanhou todo o percurso
da licenciatura.

Aos meus \textbf{colegas e amigos}, pela partilha, pela motivação e pelas
conversas que, direta ou indiretamente, contribuíram para as decisões
aqui documentadas.

\vspace{1cm}
\begin{flushright}
\textit{A todos, o meu sincero obrigado.}
\end{flushright}


\cleardoublepage
\pdfbookmark[0]{Índice}{toc}
\tableofcontents
\cleardoublepage
\listoffigures
\cleardoublepage
\listoftables
\cleardoublepage
\lstlistoflistings
\cleardoublepage
% =====================================================================
%  LISTA DE ACRÓNIMOS E ABREVIATURAS
%  (Implementada como tabela manual para compilar sem makeglossaries.
%   Se preferires o pacote glossaries, é simples migrar esta lista.)
% =====================================================================
\chapter*{Lista de Acrónimos}
\addcontentsline{toc}{chapter}{Lista de Acrónimos}
\markboth{Lista de Acrónimos}{Lista de Acrónimos}

\begin{longtable}{@{}p{2.6cm}p{12cm}@{}}
\textbf{Acrónimo} & \textbf{Significado} \\
\midrule
\endhead
API    & \textit{Application Programming Interface} \\
ARIA   & \textit{Accessible Rich Internet Applications} \\
CORS   & \textit{Cross-Origin Resource Sharing} \\
CPU    & \textit{Central Processing Unit} \\
CRUD   & \textit{Create, Read, Update, Delete} \\
CSPRNG & \textit{Cryptographically Secure Pseudo-Random Number Generator} \\
DNA    & \textit{Deoxyribonucleic Acid} \\
E2E    & \textit{End-to-End} (testes ponta-a-ponta) \\
EXIF   & \textit{Exchangeable Image File Format} \\
FFmpeg & \textit{Fast Forward MPEG} (suíte de processamento multimédia) \\
GEDCOM & \textit{Genealogical Data Communication} \\
GPS    & \textit{Global Positioning System} \\
GPU    & \textit{Graphics Processing Unit} \\
HTTP(S)& \textit{HyperText Transfer Protocol (Secure)} \\
IA     & Inteligência Artificial \\
JSON   & \textit{JavaScript Object Notation} \\
JWKS   & \textit{JSON Web Key Set} \\
JWT    & \textit{JSON Web Token} \\
LLM    & \textit{Large Language Model} (Modelo de Linguagem de Grande Escala) \\
MLLM   & \textit{Multimodal Large Language Model} \\
MP4    & \textit{MPEG-4 Part 14} (contentor de vídeo) \\
OCR    & \textit{Optical Character Recognition} \\
ORM    & \textit{Object-Relational Mapping} \\
PT-PT  & Português Europeu \\
RAG    & \textit{Retrieval-Augmented Generation} \\
RGPD   & Regulamento Geral sobre a Proteção de Dados \\
RLS    & \textit{Row-Level Security} \\
REST   & \textit{Representational State Transfer} \\
RF     & Requisito Funcional \\
RNF    & Requisito Não Funcional \\
SPA    & \textit{Single-Page Application} \\
SQL    & \textit{Structured Query Language} \\
TLS    & \textit{Transport Layer Security} \\
TTS    & \textit{Text-to-Speech} (Síntese de Voz) \\
UI/UX  & \textit{User Interface / User Experience} \\
UUID   & \textit{Universally Unique Identifier} \\
WCAG   & \textit{Web Content Accessibility Guidelines} \\
\end{longtable}


% --------------------------- CORPO ----------------------------------
\cleardoublepage
\pagenumbering{arabic}
\setcounter{page}{1}

% =====================================================================
%  CAPÍTULO 1 — INTRODUÇÃO
% =====================================================================
\chapter{Introdução}
\label{cap:introducao}

\section{Contextualização}
\label{sec:contextualizacao}

O interesse global pela genealogia e pela história familiar cresceu de
forma acentuada nas últimas décadas. A digitalização massiva de registos
históricos e a democratização do acesso à internet permitiram que
praticamente qualquer pessoa pudesse investigar as suas origens; em
paralelo, os testes de \textsc{dna} de ascendência abriram caminho para
que muitos, com origens até então desconhecidas, encontrassem finalmente
respostas. Plataformas como a \emph{Ancestry}, a \emph{MyHeritage}, a
\emph{FamilySearch} e a \emph{Findmypast} agregam hoje milhares de milhões
de registos e sustentam uma verdadeira era moderna de descoberta
familiar~\cite{familytreemag}.

Paralelamente a este movimento, a esmagadora maioria das memórias
familiares contemporâneas existe já em formato digital — fotografias no
telemóvel, vídeos no computador, documentos digitalizados. Contudo, esse
material encontra-se tipicamente \emph{disperso}, \emph{desorganizado} e
\emph{vulnerável}: vive espalhado por dispositivos, serviços de nuvem e
álbuns esquecidos, sem datação coerente, sem identificação das pessoas e,
sobretudo, sem qualquer fio narrativo que lhe dê sentido. À medida que os
membros mais velhos de uma família desaparecem, perde-se com eles o
contexto que tornava aquelas imagens significativas.

Os avanços recentes em \textbf{inteligência artificial generativa} — em
particular os \emph{Large Language Models} (\textsc{llm}) e os modelos
multimodais capazes de interpretar imagens — criam, pela primeira vez, a
possibilidade técnica de transformar este arquivo passivo num produto
narrativo vivo. É neste cruzamento entre preservação da memória familiar
e IA generativa que se enquadra o presente projeto.

\section{Definição do Problema}
\label{sec:problema}

As soluções existentes no mercado abordam apenas parcialmente o
problema. As grandes plataformas de genealogia são excelentes a
\emph{agregar registos} e a \emph{construir árvores}, mas não produzem
narrativa. Ferramentas de consumo como o \emph{Google~Photos} ou o
\emph{Apple~Memories} geram \emph{slideshows} automáticos visualmente
agradáveis, mas limitam-se a selecionar fotografias por critérios
algorítmicos, sem construir uma história coerente nem integrar dados de
múltiplas fontes (fotografias, documentos, árvores genealógicas). O que
falta — e que constitui o problema central deste trabalho — é a
\textbf{profundidade narrativa}: a capacidade de cruzar material
heterogéneo e dele extrair, de forma automática, uma história
factualmente ancorada, emocionalmente significativa e linguisticamente
correta.

Formalmente, o problema pode enunciar-se assim: \emph{dado um conjunto
heterogéneo e não estruturado de material de arquivo familiar
(fotografias com metadados, documentos digitalizados e uma árvore
genealógica), como produzir automaticamente — sem intervenção manual de
escrita — uma narrativa textual coerente e um vídeo documental narrado,
factualmente fiéis ao material fornecido e expressos em português
europeu?}

Este problema levanta vários desafios técnicos não triviais, que se
discutem em detalhe no \Cref{cap:estado_arte}: a \emph{alucinação} dos
modelos generativos (invenção de factos), a \emph{coerência temporal}
entre fontes com datas conflituosas, a \emph{identificação de pessoas} no
material e a própria \emph{avaliação da qualidade} narrativa.

\section{Objetivos}
\label{sec:objetivos}

\subsection{Objetivo geral}

Conceber, desenvolver e validar um sistema que receba material familiar
já existente e o transforme automaticamente numa narrativa multimédia —
concretamente, num vídeo documental narrado — recorrendo a inteligência
artificial generativa, com uma interface acessível e intuitiva.

\subsection{Objetivos específicos}

\begin{enumerate}[label=\textbf{O\arabic*.},leftmargin=*]
  \item \textbf{Ingerir e compreender material multimodal}: processar
  fotografias (metadados \textsc{exif} e conteúdo visual), documentos
  digitalizados (\textsc{ocr}) e árvores genealógicas (\textsc{gedcom}).
  \item \textbf{Estruturar temporalmente} os factos extraídos, resolvendo
  datas incompletas ou conflituosas e modelando explicitamente as
  relações de parentesco num grafo familiar.
  \item \textbf{Gerar narrativas} fluidas e factualmente ancoradas, em
  português europeu, controlando o tom e a estrutura através de
  \emph{templates} temáticos e de técnicas de \emph{prompt engineering}.
  \item \textbf{Produzir vídeos documentais} que combinem as fotografias
  enquadradas (com movimento opcional, efeito \emph{Ken~Burns}), a narração
  por síntese de voz e elementos cinematográficos (títulos, legendas,
  transições).
  \item \textbf{Disponibilizar} estas capacidades através de uma aplicação
  Web moderna, segura, multilingue e implantada na nuvem.
  \item \textbf{Garantir baixo custo e privacidade opcional}, assentando em
  ferramentas de código aberto e em camadas gratuitas, e mantendo a
  inferência \emph{local} como alternativa para operação privada e
  \emph{offline}.
\end{enumerate}

\section{Resultados Desejados e Âmbito}
\label{sec:ambito}

O âmbito do projeto, definido em articulação com a orientação, fixa
quatro tipos de \emph{input} — fotografias, vídeos, documentos e
informação genealógica — e um \emph{output} principal: um pequeno vídeo
documental ou \emph{slideshow} narrado, gerado automaticamente a partir
desse material. Dadas as limitações tecnológicas (e de custo) da geração
de vídeo longo por IA, o projeto concentra-se deliberadamente em
\emph{incrementos narrativos curtos} — segmentos de história que podem,
eventualmente, ser combinados numa narrativa mais longa. Esta opção é
academicamente defensável: um \emph{slideshow} narrado de alta qualidade,
construído programaticamente, cumpre integralmente os objetivos sem
depender de APIs de geração de vídeo dispendiosas.

\section{Contribuições}
\label{sec:contribuicoes}

As principais contribuições deste trabalho são:

\begin{itemize}[leftmargin=*]
  \item Uma \textbf{arquitetura modular em \emph{pipeline}} (M1--M4) que
  separa claramente ingestão, organização temporal, geração textual e
  produção multimédia.
  \item Um \textbf{subsistema RAG multi-inquilino com degradação
  graciosa}, que mantém o sistema funcional mesmo sem base vetorial.
  \item Um conjunto de \textbf{\emph{templates} de geração calibrados para
  português europeu}, que mitigam a tendência dos modelos para o
  português do Brasil.
  \item Uma \textbf{estratégia de IA em cascata e a custo nulo} (Groq na
  nuvem, Ollama local e Gemini), que concilia desempenho, robustez por
  \emph{fallback} e privacidade opcional (inferência local).
  \item Uma \textbf{narrativa segmentada em cenas} que liga cada parágrafo
  às fotografias que o ilustram, permitindo sincronizar imagem e narração
  no vídeo — a abordagem que distingue um documentário de um simples
  \emph{slideshow}.
  \item Uma \textbf{aplicação Web completa e implantada}, do
  \emph{frontend} ao \emph{deployment} na nuvem, com \textbf{geração
  assíncrona} (Celery ou executor \emph{in-process}) e resiliência a
  \emph{cold starts}.
\end{itemize}

\section{Estrutura do Documento}
\label{sec:estrutura}

O presente relatório encontra-se organizado em oito capítulos. O
\textbf{\Cref{cap:introducao}} (este) apresenta o contexto, o problema e
os objetivos. O \textbf{\Cref{cap:estado_arte}} analisa as soluções
existentes e enquadra teoricamente os conceitos fundamentais (LLM,
multimodalidade, RAG, geração de narrativa e de vídeo). O
\textbf{\Cref{cap:metodologia}} descreve a metodologia de desenvolvimento
e justifica a \emph{stack} tecnológica. O \textbf{\Cref{cap:requisitos}}
formaliza os requisitos funcionais e não funcionais e os casos de uso. O
\textbf{\Cref{cap:arquitetura}} apresenta a arquitetura do sistema, o
modelo de dados e o \emph{design} de interface. O
\textbf{\Cref{cap:implementacao}} detalha a implementação dos quatro
módulos e as decisões técnicas, com excertos de código críticos. O
\textbf{\Cref{cap:testes}} cobre a validação e os testes. Por fim, o
\textbf{\Cref{cap:conclusoes}} sintetiza os resultados, discute as
limitações e aponta o trabalho futuro.

% =====================================================================
%  CAPÍTULO 2 — ESTADO DA ARTE / REVISÃO BIBLIOGRÁFICA
% =====================================================================
\chapter{Estado da Arte e Enquadramento Teórico}
\label{cap:estado_arte}

Este capítulo divide-se em duas partes. A primeira (\Cref{sec:solucoes})
analisa as soluções existentes no mercado que se aproximam, total ou
parcialmente, dos objetivos deste projeto, identificando o que fazem bem
e o que lhes falta. A segunda (\Cref{sec:enquadramento}) introduz os
conceitos teóricos fundamentais — modelos de linguagem, multimodalidade,
geração de narrativa e geração de vídeo — necessários para compreender as
opções técnicas tomadas nos capítulos seguintes.

\section{Análise de Soluções Existentes}
\label{sec:solucoes}

\subsection{Os ``Big Four'' da genealogia digital}
\label{sub:bigfour}

O mercado da genealogia digital é dominado por quatro grandes plataformas
que se distinguem pela escala dos seus arquivos de registos históricos,
pela extensão das árvores genealógicas colaborativas e pelas ligações
genéticas que oferecem~\cite{familytreemag}:

\begin{itemize}[leftmargin=*]
  \item \textbf{Ancestry.com} — dominante no mercado norte-americano,
  forte integração com testes de \textsc{dna}, particularmente eficaz na
  investigação de ascendência nos EUA.
  \item \textbf{MyHeritage} — mais vocacionada para famílias
  internacionais, reconhecida pelas suas ferramentas de fotografia e
  \textsc{dna}.
  \item \textbf{FamilySearch} — completamente gratuita, mantida pela
  Igreja SUD, com um alcance verdadeiramente global.
  \item \textbf{Findmypast} — especializada em registos do Reino Unido e
  da Irlanda.
\end{itemize}

Estas plataformas são exemplares na agregação e indexação de registos,
mas o seu foco é a \emph{descoberta} de factos genealógicos, não a
\emph{construção de narrativa}.

\subsection{MyHeritage — o caso de estudo mais relevante}
\label{sub:myheritage}

De entre as quatro, a \textbf{MyHeritage} é a plataforma mais próxima
daquilo que este projeto pretende construir, sobretudo pelo seu
investimento em funcionalidades de IA generativa aplicada à memória
familiar~\cite{myheritage_rootstech}:

\begin{itemize}[leftmargin=*]
  \item \textbf{Deep Nostalgia} (2020) — anima fotografias antigas,
  conferindo movimento facial realista a retratos estáticos.
  \item \textbf{DeepStory} (2022) — permite que os antepassados
  ``contem'' as suas próprias histórias de vida, gerando vídeo a partir
  de retratos e de texto.
  \item \textbf{AI Biographer} (2024) — gera, para qualquer antepassado,
  uma entrada de estilo enciclopédico, com contexto histórico,
  fotografias e notas de rodapé que citam a origem dos factos
  genealógicos. Nos dois primeiros meses após o lançamento foram geradas
  cerca de 100.000 biografias~\cite{myheritage_rootstech}.
\end{itemize}

O \emph{AI Biographer} é particularmente relevante por demonstrar a
viabilidade comercial da geração automática de texto biográfico
\emph{ancorado em factos}. Distingue-se, contudo, do presente projeto em
dois aspetos: foca-se em texto (e não em vídeo documental narrado a
partir do arquivo do utilizador) e opera dentro de um ecossistema fechado.

\subsection{Google Photos e Apple Memories — o outro lado}
\label{sub:googlephotos}

Do lado das aplicações de consumo, a IA integrada em serviços como o
\emph{Google~Photos} reanima automaticamente o passado, selecionando
algoritmicamente fotografias que evocam pessoas, locais e momentos. A
funcionalidade \emph{Cinematic Photos and Memories} recorre a
\emph{machine learning} para gerar pequenos clipes animados com um efeito
de profundidade 3D; segundo a Google, o envolvimento dos utilizadores com
a funcionalidade \emph{Memories} cresceu 40\% num ano. Estas soluções são
tecnicamente sofisticadas e visualmente agradáveis, mas o seu resultado é
um \emph{slideshow} estético, não uma \emph{história}.

\subsection{Síntese comparativa e lacuna identificada}
\label{sub:lacuna}

A \Cref{tab:comparativo} sintetiza a análise. O denominador comum às
soluções existentes é a ausência de \textbf{profundidade narrativa} e de
\textbf{integração multifonte}: nenhuma cruza simultaneamente
fotografias, documentos digitalizados, árvores genealógicas e metadados
para construir uma narrativa coerente e factualmente ancorada. É
precisamente essa lacuna que este projeto se propõe preencher.

\begin{table}[H]
\centering
\caption{Análise comparativa das soluções existentes face aos objetivos
do projeto.}
\label{tab:comparativo}
\small
\begin{tabularx}{\linewidth}{@{}p{3.1cm}YYcc@{}}
\toprule
\textbf{Solução} & \textbf{Foco principal} & \textbf{Integra
multifonte?} & \textbf{Narrativa} & \textbf{Vídeo} \\
\midrule
Ancestry / FamilySearch / Findmypast & Registos e árvores genealógicas &
Parcial & Não & Não \\
MyHeritage (AI Biographer) & Biografia textual ancorada & Sim
(genealogia) & Sim (texto) & Parcial \\
Google Photos / Apple Memories & \emph{Slideshows} automáticos & Não &
Não & Sim (estético) \\
\textbf{Este projeto} & Narrativa multimédia do arquivo & \textbf{Sim} &
\textbf{Sim} & \textbf{Sim (documental)} \\
\bottomrule
\end{tabularx}
\end{table}

\section{Enquadramento Teórico}
\label{sec:enquadramento}

\subsection{Modelos de Linguagem de Grande Escala (LLM)}
\label{sub:llm}

Um \emph{Large Language Model} (\textsc{llm}) é, na sua essência, um
modelo treinado para prever a palavra (ou \emph{token}) seguinte numa
sequência de texto. Treinado sobre quantidades massivas de texto —
páginas da internet, livros, artigos científicos — o modelo aprende
padrões estatísticos da linguagem ao ponto de conseguir gerar texto novo,
coerente e contextualmente adequado~\cite{brown2020gpt3}. O mecanismo
central que tornou estes modelos viáveis é a \emph{self-attention},
introduzida na arquitetura \emph{Transformer}~\cite{vaswani2017attention},
que permite ao modelo ponderar simultaneamente as relações entre todas as
palavras de uma sequência, atribuindo pesos diferentes a cada par.

A forma de comunicar com um \textsc{llm} é o \emph{prompt} — a instrução
fornecida ao modelo. A qualidade do resultado depende fortemente da
qualidade do \emph{prompt}, uma constatação que deu origem a uma área de
investigação própria, o \textbf{\emph{prompt engineering}}. A diferença é
ilustrativa: um \emph{prompt} fraco como \emph{``escreve sobre uma
família''} produz texto genérico; um \emph{prompt} bem construído, que
forneça factos concretos, defina a perspetiva, o tom, o formato e a
estrutura, produz resultados radicalmente superiores. Esta observação é
central para o \Cref{cap:implementacao}, onde se descreve o sistema de
\emph{templates} do módulo de geração narrativa.

\subsection{Modelos multimodais — ver além do texto}
\label{sub:multimodal}

Os modelos modernos já não processam apenas texto. Os \emph{Multimodal
Large Language Models} (\textsc{mllm}) recebem como entrada texto, imagem
e até áudio e vídeo em simultâneo. Quando uma fotografia é apresentada a
um modelo como o \emph{GPT-4V} ou o \emph{Gemini}, ocorre, em termos
simplificados, o seguinte~\cite{yin2023mllmsurvey}:

\begin{enumerate}[leftmargin=*]
  \item A imagem é dividida em pequenos blocos (\emph{patches}).
  \item Cada \emph{patch} é convertido numa representação numérica
  (\emph{embedding} visual).
  \item Esses \emph{embeddings} visuais são projetados para o mesmo
  espaço dos \emph{embeddings} de texto.
  \item O modelo processa texto e imagem conjuntamente, como se fossem a
  mesma modalidade.
\end{enumerate}

Representado por modelos como o \emph{GPT-4V}, o paradigma \textsc{mllm}
tornou-se um foco crescente de investigação, usando \textsc{llm}
poderosos como ``cérebro'' para executar tarefas multimodais. Capacidades
emergentes surpreendentes — como \emph{escrever histórias a partir de
imagens} — sugerem um caminho promissor~\cite{yin2023mllmsurvey}. É
exatamente esta capacidade que o módulo de ingestão (M1) explora: um
modelo de visão descreve cada fotografia, gerando matéria-prima factual
para a posterior geração narrativa.

\subsection{Geração automática de narrativas}
\label{sub:geracao_narrativa}

Uma narrativa não é uma mera sequência de factos. Uma boa história possui
\emph{estrutura} (início, desenvolvimento, clímax, conclusão),
\emph{perspetiva} (quem narra e de que ponto de vista), \emph{coerência}
(os factos não se contradizem), \emph{arco emocional} (tensão e resolução)
e \emph{contexto} (os eventos ligam-se a um mundo mais amplo). O desafio
da geração automática consiste em obter tudo isto a partir de dados
brutos e desorganizados.

Identificam-se três paradigmas principais de geração de narrativas:

\begin{description}[leftmargin=*]
  \item[Paradigma 1 — \emph{Template-Based}.] O sistema mais simples:
  definem-se modelos fixos e preenchem-se os espaços com dados. Para este
  projeto, os \emph{templates} são o ponto de partida da arquitetura —
  definem a estrutura narrativa (cronológica, por personagem, temática) e
  o \textsc{llm} preenche-a com linguagem natural rica.
  \item[Paradigma 2 — \emph{Retrieval-Augmented Generation} (RAG).]
  Combina uma fase de \emph{retrieval} (busca de informação relevante numa
  base de dados de factos) com uma fase de \emph{generation} (o
  \textsc{llm} usa essa informação para gerar texto)~\cite{lewis2020rag}.
  \item[Paradigma 3 — \emph{Fine-tuning} e modelos especializados.]
  Parte-se de um modelo base (por exemplo o LLaMA~\cite{touvron2023llama})
  e treina-se especificamente em histórias familiares, para que ``fale''
  com o estilo de um biógrafo. É a abordagem mais poderosa, mas também a
  mais exigente em dados e computação.
\end{description}

Para o âmbito deste projeto, o \textbf{RAG é a abordagem mais
pragmática}: não exige treino de modelos próprios, funciona bem tanto com
modelos locais como com APIs existentes, e — crucialmente — ancora a
geração em factos verificados, mitigando a alucinação. O
\Cref{sub:desafios} aprofunda esta justificação.

\subsection{Desafios técnicos da geração narrativa}
\label{sub:desafios}

\begin{description}[leftmargin=*]
  \item[Alucinação.] Os \textsc{llm} podem inventar factos inexistentes;
  na ausência de informação sobre um período, tendem a ``preencher'' com
  eventos plausíveis mas falsos. \emph{Mitigação:} RAG com factos
  verificados — instruir o modelo a usar \emph{apenas} a informação
  fornecida e a admitir a ausência de registo.
  \item[Coerência temporal.] Dados de décadas podem conter
  inconsistências (uma pessoa que aparece numa fotografia datada antes do
  seu nascimento). \emph{Mitigação:} um módulo de validação que verifica
  a consistência da linha temporal antes da geração (papel do M2).
  \item[Identificação de pessoas.] Sem saber quem é quem nas fotografias,
  a narrativa torna-se genérica. \emph{Mitigação:} reconhecimento facial
  (e.g.\ \emph{DeepFace}) combinado com confirmação manual do utilizador.
  Esta funcionalidade situa-se fora do âmbito da versão atual e é
  discutida como trabalho futuro (\Cref{cap:conclusoes}).
  \item[Avaliação de qualidade.] Determinar se uma narrativa gerada é
  ``boa'' exige métricas de avaliação — objetivas (e.g.\ aderência
  linguística) e subjetivas (painéis humanos) — discutidas no
  \Cref{cap:testes}.
\end{description}

\subsection{Geração de vídeo: pipeline programático vs.\ IA de vídeo}
\label{sub:geracao_video}

Existem duas abordagens fundamentalmente distintas para transformar
fotografias num vídeo:

\begin{description}[leftmargin=*]
  \item[Abordagem A — \emph{pipeline} programático.] Monta-se o vídeo
  como faria um editor: define-se a duração de cada fotografia,
  adiciona-se narração por \textsc{tts}, música de fundo e o efeito
  \emph{Ken~Burns} (zoom/panorâmica lentos que dão vida a imagens
  estáticas), exportando um ficheiro \textsc{mp4}. É o caminho com
  bibliotecas de código aberto como o \emph{MoviePy}~\cite{moviepy} e o
  \emph{FFmpeg}~\cite{ffmpeg}, de custo nulo e plenamente controlável.
  \item[Abordagem B — geração de vídeo por IA.] Recorre a APIs como o
  \emph{Runway}~\cite{runway}, \emph{Pika}, \emph{Luma} ou \emph{Kling} para, a partir
  de uma fotografia estática, gerar segundos de vídeo animado com
  movimento realista. É mais impressionante, mas também mais cara e
  complexa, e levanta questões de consistência e de custo por geração.
\end{description}

Para um projeto académico, um \emph{slideshow} narrado de qualidade
construído pela Abordagem A é \textbf{suficiente e academicamente
defensável}, cumprindo os objetivos sem dependência de APIs dispendiosas.
Foi essa a opção adotada, conforme detalhado no \Cref{cap:implementacao}.
A \Cref{tab:videoapis} resume, para referência, as principais APIs de
geração de vídeo da Abordagem B.

\begin{table}[H]
\centering
\caption{Principais APIs de geração de vídeo por IA (Abordagem B), para
referência. Não utilizadas na versão atual por motivos de custo.}
\label{tab:videoapis}
\small
\begin{tabularx}{\linewidth}{@{}lYl@{}}
\toprule
\textbf{Ferramenta} & \textbf{Ideal para} & \textbf{API} \\
\midrule
Pika Labs            & Prototipagem rápida                  & Sim \\
Runway Gen-4         & Qualidade profissional               & Sim \\
Luma (Dream Machine) & Conversão foto $\rightarrow$ vídeo realista & Sim \\
Kling AI             & Alternativa de boa qualidade         & Sim \\
\bottomrule
\end{tabularx}
\end{table}

\subsection{Síntese de voz (Text-to-Speech)}
\label{sub:tts}

A síntese de voz (\textsc{tts}) converte o texto narrativo em fala,
dando ``voz'' à história. A qualidade da voz afeta diretamente o impacto
emocional do vídeo: uma narração natural e expressiva tem um efeito
completamente distinto de uma voz robótica. Entre as opções de código
aberto/gratuitas consideradas: o \emph{gTTS} (simples, dependente de
internet, qualidade razoável mas pouco expressiva), as vozes neuronais do
\emph{edge-tts} (qualidade elevada, vozes europeias) e o \emph{Coqui~TTS}
(local, com a possibilidade de \emph{voice cloning} a partir de uma
amostra real — funcionalidade tecnicamente e eticamente rica, mas
exigente em recursos). As opções efetivamente integradas no projeto
discutem-se no \Cref{cap:implementacao}.

\subsection{Ferramentas de código aberto por módulo}
\label{sub:ferramentas}

Um princípio orientador do projeto é o custo essencialmente nulo, através
de ferramentas de código aberto e camadas gratuitas. A
\Cref{tab:ferramentas} resume, por módulo, as tecnologias consideradas no
estudo inicial — várias delas efetivamente adotadas.

\begin{table}[H]
\centering
\caption{Ferramentas de código aberto / gratuitas consideradas, por
módulo.}
\label{tab:ferramentas}
\small
\begin{tabularx}{\linewidth}{@{}p{3.4cm}Y@{}}
\toprule
\textbf{Módulo} & \textbf{Ferramentas} \\
\midrule
M1 — Ingestão & Google Gemini (visão, camada gratuita); Llama 3.2 Vision
via Ollama (local); extração \textsc{exif}; Tesseract \textsc{ocr};
\emph{parser} GEDCOM. \\
M2 — Temporal & Lógica própria em Python; NetworkX para o grafo de
parentesco. \\
M3 — Narrativa & Groq (Llama 3.3 70B, camada gratuita na nuvem); Ollama +
Llama 3.2 / Mistral / Gemma (local); Google Gemini Flash (camada gratuita);
Hugging Face Inference (com limites). \\
M4 — Multimédia & MoviePy; gTTS / edge-tts / Coqui TTS; FFmpeg. \\
\bottomrule
\end{tabularx}
\end{table}

\noindent Onde poderiam surgir custos — APIs comerciais de \textsc{llm}
proprietárias em produção, ou APIs de geração de vídeo (Runway, Pika) —
o projeto evita-os deliberadamente: modelos de camada gratuita (\emph{Groq}
na nuvem e \emph{Ollama} em local) para o texto e um \emph{slideshow}
programático com MoviePy no lugar da geração de vídeo por IA.

% =====================================================================
%  CAPÍTULO 3 — METODOLOGIA E TECNOLOGIAS
% =====================================================================
\chapter{Metodologia e Tecnologias}
\label{cap:metodologia}

\section{Metodologia de Desenvolvimento}
\label{sec:metodologia}

O desenvolvimento seguiu uma abordagem \textbf{ágil e incremental},
adaptada às características de um projeto académico individual. Em vez de
um modelo em cascata rígido, adotou-se um ciclo iterativo inspirado em
\emph{Kanban}: o trabalho foi organizado em incrementos verticais, cada
um entregando uma fatia funcional ponta-a-ponta (da ingestão à geração),
o que permitiu validar continuamente o \emph{pipeline} completo em vez de
esperar pela conclusão de todos os módulos.

\subsection{Fases do projeto}

O projeto desenrolou-se em fases sobrepostas:

\begin{enumerate}[label=\textbf{F\arabic*.},leftmargin=*]
  \item \textbf{Investigação e estado da arte} — levantamento das
  soluções existentes (\Cref{cap:estado_arte}), das tecnologias de IA
  generativa e das ferramentas de código aberto por módulo.
  \item \textbf{Prototipagem do \emph{pipeline}} — implementação inicial
  dos quatro módulos (M1--M4) em modo \emph{local-first}, com SQLite e
  ficheiros locais, validando a viabilidade técnica end-to-end.
  \item \textbf{Consolidação do \emph{backend}} — introdução de
  autenticação, fila de tarefas assíncronas, validação de \emph{uploads},
  \emph{rate limiting}, \emph{logging} estruturado e testes automatizados.
  \item \textbf{Construção do \emph{frontend}} — desenvolvimento da
  aplicação React completa (páginas, gestão de estado,
  internacionalização, tema).
  \item \textbf{Migração para a nuvem} — evolução para uma arquitetura
  multi-inquilino com Supabase (autenticação, base de dados e
  armazenamento) e \emph{deployment} em Render e Vercel.
  \item \textbf{Validação e documentação} — testes, recolha de resultados
  e redação do presente relatório.
\end{enumerate}

\subsection{Gestão e controlo de versões}

Todo o código foi versionado com \textbf{Git}, com \emph{commits}
descritivos a marcar cada incremento funcional e cada correção. O
histórico de versões serviu simultaneamente de registo de progresso e de
mecanismo de segurança (possibilidade de reverter alterações). A
\Cref{fig:metodologia} ilustra o cronograma das fases.

\begin{figure}[H]
\centering
% \includegraphics[width=\linewidth]{cronograma.png}
\figph[5.5cm]{Cronograma do projeto (diagrama de Gantt) com as fases
F1--F6 ao longo do tempo. Sugestão: criar em ferramenta de gestão de
projeto ou em TikZ.}
\caption{Cronograma das fases de desenvolvimento.}
\label{fig:metodologia}
\end{figure}

\section{\emph{Stack} Tecnológica e Justificação}
\label{sec:stack}

A seleção tecnológica obedeceu a três critérios: \emph{(i)} maturidade e
adequação técnica; \emph{(ii)} custo nulo ou camada gratuita; e
\emph{(iii)} possibilidade de execução local para garantir privacidade.
A \Cref{tab:stack} resume a \emph{stack}; as subsecções seguintes
justificam as escolhas principais.

\begin{table}[H]
\centering
\caption{\emph{Stack} tecnológica por camada.}
\label{tab:stack}
\small
\begin{tabularx}{\linewidth}{@{}p{3cm}Y@{}}
\toprule
\textbf{Camada} & \textbf{Tecnologias} \\
\midrule
\emph{Frontend} & React~18, TypeScript (modo estrito), Vite~5, Tailwind
CSS~3, TanStack Query, Zustand, React Router~6, react-i18next, Framer
Motion, React Flow, Sonner, Lucide. \\
\emph{Backend} & Python~3, FastAPI (assíncrono), SQLAlchemy~2 (async),
Pydantic~v2, slowapi, structlog, Celery + Redis (opcional). \\
IA / ML & Ollama (Llama~3.2~3B, local); Google Gemini~1.5~Flash (nuvem);
ChromaDB + \emph{sentence-transformers} (RAG, opcional); Tesseract
(\textsc{ocr}); NetworkX (grafo). \\
Multimédia & MoviePy + FFmpeg; Pillow (Ken Burns / \emph{letterbox});
edge-tts / gTTS (\textsc{tts}). \\
Dados & Supabase Auth (JWT/JWKS), PostgreSQL (asyncpg), Supabase Storage;
SQLite (desenvolvimento local). \\
Operação & Docker; Render (\emph{backend}); Vercel (\emph{frontend});
Supabase (dados); Git. \\
\bottomrule
\end{tabularx}
\end{table}

\subsection{\emph{Backend}: FastAPI}

Escolheu-se o \textbf{FastAPI} pela sua natureza \emph{assíncrona}
(essencial para um sistema com operações de I/O e de IA potencialmente
longas), pela validação automática de dados via \textbf{Pydantic} (que
elimina uma classe inteira de erros de \emph{input}) e pela geração
automática de documentação interativa (\emph{Swagger UI}). A combinação
com \textbf{SQLAlchemy} assíncrono e \textbf{Pydantic~v2} oferece tipagem
forte de ponta a ponta no servidor.

\subsection{IA: estratégia híbrida local/nuvem}

A decisão mais estruturante foi a \textbf{estratégia de IA híbrida}.
Privilegia-se a inferência \textbf{local} via \textbf{Ollama} (executando
o \emph{Llama~3.2~3B}) por três razões: \emph{privacidade} (os factos do
arquivo não saem da máquina), \emph{custo} (nulo) e \emph{autonomia} (não
depende de conectividade). Como \emph{fallback}, recorre-se ao
\textbf{Google Gemini~1.5~Flash}, cuja camada gratuita é suficiente para
um protótipo académico e cuja maior capacidade é útil para narrativas
mais exigentes ou quando não há \textsc{gpu} local. Para a análise visual
das fotografias (M1) usa-se o \textbf{Gemini Vision}. Esta arquitetura
de dupla via é detalhada no \Cref{cap:implementacao}.

\subsection{RAG: ChromaDB}

Para o subsistema \textbf{RAG} adotou-se o \textbf{ChromaDB} como base
vetorial, com \emph{embeddings} de \emph{sentence-transformers}. A
escolha recaiu sobre o ChromaDB pela simplicidade de integração
(persistência local em ficheiro, sem servidor dedicado). Como se verá, o
subsistema foi desenhado para \emph{degradar graciosamente} na ausência
desta dependência, uma decisão motivada pelas restrições de
\emph{deployment} discutidas no \Cref{cap:implementacao}.

\subsection{Multimédia: MoviePy, FFmpeg e Pillow}

Para a Abordagem A (\Cref{sub:geracao_video}) usaram-se o \textbf{MoviePy}
(composição programática de vídeo), o \textbf{FFmpeg} (codificação e
mistura de áudio) e o \textbf{Pillow} (processamento de imagem para o
efeito \emph{Ken~Burns} e para o enquadramento \emph{letterbox}). Para a
narração, optou-se pelas vozes neuronais europeias do \textbf{edge-tts},
com o \textbf{gTTS} como alternativa.

\subsection{\emph{Frontend}: React, TypeScript e TanStack Query}

No \emph{frontend} escolheu-se o \textbf{React} com \textbf{TypeScript em
modo estrito} para capturar, em tempo de compilação, incompatibilidades
de esquema com o \emph{backend}. O \textbf{Vite} foi preferido como
\emph{bundler} pela velocidade do \emph{Hot Module Replacement}. Para a
gestão de estado, distinguiu-se rigorosamente o \emph{estado de servidor}
— gerido pelo \textbf{TanStack Query} (cache, invalidação,
\emph{polling}) — do \emph{estado de cliente} — gerido pelo
\textbf{Zustand} (autenticação e tema). Esta separação, hoje considerada
boa prática, evita a complexidade de soluções como o Redux para dados que
são, na sua maioria, \emph{server-state}.

\subsection{Dados e operação: Supabase, Docker, Render e Vercel}

A camada de dados foi externalizada para o \textbf{Supabase}, que fornece
de forma integrada autenticação (com tokens JWT validados por JWKS),
base de dados \textbf{PostgreSQL} e armazenamento de objetos
(\emph{Storage}). O \emph{backend} é empacotado em \textbf{Docker} (para
incluir dependências de sistema como FFmpeg, Tesseract e libmagic) e
implantado no \textbf{Render}; o \emph{frontend} é servido pela
\textbf{Vercel}. Esta combinação assenta integralmente em camadas
gratuitas, mantendo o custo operacional próximo de zero.

% =====================================================================
%  CAPÍTULO 4 — LEVANTAMENTO E ANÁLISE DE REQUISITOS
% =====================================================================
\chapter{Levantamento e Análise de Requisitos}
\label{cap:requisitos}

Este capítulo formaliza os requisitos do sistema, distinguindo os
\emph{funcionais} (o que o sistema deve fazer) dos \emph{não funcionais}
(restrições de qualidade), e descreve os principais casos de uso através
de um diagrama e de \emph{user stories}.

\section{Requisitos Funcionais}
\label{sec:rf}

Os requisitos funcionais derivam diretamente dos objetivos
(\Cref{sec:objetivos}) e do âmbito definido com a orientação. A
\Cref{tab:rf} lista-os, atribuindo a cada um uma prioridade
(\emph{Must}, \emph{Should} ou \emph{Could}, segundo a notação MoSCoW).

\begin{longtable}{@{}p{1.4cm}p{8.8cm}p{2.4cm}@{}}
\caption{Requisitos funcionais do sistema.}\label{tab:rf}\\
\toprule
\textbf{ID} & \textbf{Descrição} & \textbf{Prioridade} \\
\midrule
\endfirsthead
\multicolumn{3}{c}{\tablename\ \thetable\ -- continuação}\\
\toprule
\textbf{ID} & \textbf{Descrição} & \textbf{Prioridade} \\
\midrule
\endhead
\rf{1}  & O sistema deve permitir registo, início e fim de sessão, e
recuperação/alteração de palavra-passe. & Must \\
\rf{2}  & Cada utilizador deve aceder exclusivamente aos seus próprios
dados (isolamento multi-inquilino). & Must \\
\rf{3}  & O sistema deve permitir o carregamento de fotografias,
individualmente e em lote, com \emph{drag-and-drop}. & Must \\
\rf{4}  & O sistema deve importar árvores genealógicas em formato GEDCOM e
construir a estrutura familiar. & Must \\
\rf{5}  & O sistema deve extrair automaticamente metadados \textsc{exif}
das fotografias (data, GPS, câmara). & Must \\
\rf{6}  & O sistema deve descrever o conteúdo visual de cada fotografia
através de um modelo de IA de visão. & Must \\
\rf{7}  & O sistema deve extrair texto de documentos digitalizados por
\textsc{ocr}. & Should \\
\rf{8}  & O sistema deve construir uma linha temporal cronológica,
resolvendo datas incompletas. & Should \\
\rf{9}  & O sistema deve gerar narrativas a partir dos factos, permitindo
escolher tema, tom e idioma (PT/EN). & Must \\
\rf{10} & O sistema deve gerar um vídeo documental narrado a partir de uma
narrativa. & Must \\
\rf{11} & O sistema deve permitir organizar fotografias, histórias e
vídeos em projetos (coleções curadas). & Should \\
\rf{12} & O sistema deve executar gerações longas em segundo plano, com
acompanhamento de estado e notificações. & Should \\
\rf{13} & O sistema deve permitir descarregar e reproduzir os vídeos
gerados. & Must \\
\rf{14} & A interface deve estar disponível em português europeu
(omissão) e inglês. & Could \\
\rf{15} & O sistema deve expor o estado de saúde das suas dependências
(sondas de \emph{health}). & Could \\
\rf{16} & O sistema deve permitir gerir as histórias geradas --- editar,
regenerar com \emph{feedback}, marcar como favoritas, pesquisar/filtrar e
exportar para PDF. & Should \\
\rf{17} & O sistema deve permitir visualizar e editar a árvore familiar de
forma interativa (pesquisar pessoas, destacar linhagem, adicionar
familiares e exportar a árvore como imagem). & Should \\
\rf{18} & O sistema deve permitir escolher a voz da narração e incluir uma
faixa de legendas sincronizada e comutável no vídeo. & Could \\
\rf{19} & O sistema deve permitir etiquetar as pessoas presentes em cada
fotografia e gerir a galeria de cada pessoa. & Should \\
\bottomrule
\end{longtable}

\section{Requisitos Não Funcionais}
\label{sec:rnf}

\begin{longtable}{@{}p{1.6cm}p{3.4cm}p{8.6cm}@{}}
\caption{Requisitos não funcionais do sistema.}\label{tab:rnf}\\
\toprule
\textbf{ID} & \textbf{Atributo} & \textbf{Descrição} \\
\midrule
\endfirsthead
\multicolumn{3}{c}{\tablename\ \thetable\ -- continuação}\\
\toprule
\textbf{ID} & \textbf{Atributo} & \textbf{Descrição} \\
\midrule
\endhead
\rnf{1} & Segurança & Validação criptográfica de tokens; isolamento estrito
entre inquilinos; proteção contra abuso por \emph{rate limiting};
validação de \emph{uploads}. \\
\rnf{2} & Privacidade & Possibilidade de inferência local, sem envio de
dados do arquivo para serviços de terceiros; conformidade com os
princípios do RGPD. \\
\rnf{3} & Robustez & Degradação graciosa na ausência de dependências
opcionais (base vetorial, \textsc{llm} local, fila de tarefas). \\
\rnf{4} & Observabilidade & \emph{Logging} estruturado e sondas de saúde
profundas sobre todas as dependências. \\
\rnf{5} & Desempenho percebido & \emph{Cache} de estado de servidor,
\emph{polling} controlado e \emph{feedback} imediato ao utilizador. \\
\rnf{6} & Manutenibilidade & Tipagem estática (Python e TypeScript),
modularidade e cobertura de testes automatizados. \\
\rnf{7} & Portabilidade & Empacotamento em contentor Docker com todas as
dependências de sistema. \\
\rnf{8} & Usabilidade & Interface intuitiva, responsiva (\emph{desktop} e
móvel), com tema claro/escuro e suporte a acessibilidade. \\
\rnf{9} & Custo & Operação assente em ferramentas de código aberto e
camadas gratuitas (custo $\approx 0$). \\
\bottomrule
\end{longtable}

\section{Atores e Casos de Uso}
\label{sec:casos_uso}

O sistema tem um ator principal — o \textbf{Utilizador autenticado}
(o ``dono'' de um arquivo familiar) — e dois atores secundários, que são
sistemas externos: o \textbf{Serviço de IA} (Groq/Ollama/Gemini) e o
\textbf{Supabase} (autenticação, base de dados e armazenamento). A
\Cref{fig:casos_uso} apresenta o diagrama de casos de uso.

\begin{figure}[H]
\centering
% \includegraphics[width=0.92\linewidth]{casos_uso.png}
\figph[7.5cm]{Diagrama UML de casos de uso. Ator ``Utilizador''
ligado aos casos: Autenticar-se; Carregar fotografias; Importar GEDCOM;
Gerar narrativa; Gerar vídeo; Gerir projetos; Consultar tarefas; Gerir
definições. Atores secundários ``Serviço de IA'' e ``Supabase'' ligados
aos casos de geração e de persistência. Sugestão: desenhar em
draw.io / PlantUML.}
\caption{Diagrama de casos de uso do sistema.}
\label{fig:casos_uso}
\end{figure}

\section{\emph{User Stories} Principais}
\label{sec:user_stories}

As \emph{user stories} seguintes capturam a intenção do utilizador e
ligam-se aos requisitos funcionais correspondentes:

\begin{itemize}[leftmargin=*]
  \item \textbf{US1} — \emph{``Como utilizador, quero carregar as
  fotografias antigas da minha família arrastando-as para o ecrã, para
  não ter de as enviar uma a uma.''} (\rf{3})
  \item \textbf{US2} — \emph{``Como utilizador, quero importar a árvore
  genealógica que já tenho em GEDCOM, para que as narrativas conheçam as
  relações familiares.''} (\rf{4})
  \item \textbf{US3} — \emph{``Como utilizador, quero pedir uma história
  sobre um tema específico (por exemplo, a emigração do meu avô),
  escolhendo o tom, para obter uma narrativa fiel ao que tenho em
  mente.''} (\rf{9})
  \item \textbf{US4} — \emph{``Como utilizador, quero transformar uma
  história num vídeo documental narrado, para a poder partilhar e
  preservar.''} (\rf{10})
  \item \textbf{US5} — \emph{``Como utilizador, quero organizar o meu
  material em projetos temáticos, para separar diferentes ramos ou
  ocasiões da família.''} (\rf{11})
  \item \textbf{US6} — \emph{``Como utilizador, quero que as gerações
  longas corram em segundo plano e me avisem quando terminam, para não
  ficar à espera.''} (\rf{12})
  \item \textbf{US7} — \emph{``Como utilizador preocupado com a
  privacidade, quero que os meus dados não saiam da máquina sempre que
  possível.''} (\rnf{2})
\end{itemize}

% =====================================================================
%  CAPÍTULO 5 — ARQUITETURA E DESIGN DA SOLUÇÃO
% =====================================================================
\chapter{Arquitetura e Design da Solução}
\label{cap:arquitetura}

\section{Visão Global da Arquitetura}
\label{sec:visao_global}

O sistema segue uma arquitetura \textbf{cliente-servidor} de três camadas
com um \emph{pipeline} de processamento de quatro estágios no núcleo do
servidor. O cliente é uma \emph{Single-Page Application} (\textsc{spa})
React; o servidor é uma API \textsc{rest} assíncrona em FastAPI; e a
camada de dados é externalizada para serviços geridos do Supabase
(autenticação, base de dados relacional e armazenamento de objetos). A
\Cref{fig:arquitetura} (Anexo~\ref{ap:diagramas}) representa a organização
global e os fluxos principais.

Uma característica determinante da arquitetura é o modo como se gerem as
operações demoradas dos módulos M1--M4 (\textsc{ocr}, análise visual,
inferência do \textsc{llm} e montagem de vídeo), que podem demorar de
segundos a minutos. O sistema dispõe de uma \textbf{infraestrutura de
tarefas em segundo plano} --- com duas vias intermutáveis (ver
\Cref{sub:async}): uma \textbf{fila de tarefas} (Celery + Redis) quando
existe um \emph{worker} dedicado, ou um \textbf{executor \emph{in-process}}
dentro do próprio \emph{backend} quando não existe (o caso do plano
gratuito) ---, na qual o utilizador recebe de imediato um identificador de
tarefa e o \emph{frontend} acompanha o progresso por \emph{polling}. Esta
via é usada, em produção, para a \textbf{análise visual diferida das
fotografias} (M1). Em contrapartida, a \textbf{geração de narrativa e de
vídeo corre de forma síncrona}, mantendo a instância acordada durante o
processamento: uma escolha deliberada, motivada pela fragilidade dos
\emph{jobs} de segundo plano no plano gratuito (instâncias que adormecem ou
reiniciam, abandonando a tarefa --- \Cref{sub:async}). A
\Cref{fig:fullstack} (Anexo~\ref{ap:diagramas}) apresenta uma vista
detalhada, camada a camada, de toda a \emph{stack}.

\section{Evolução Arquitetural: de \emph{local-first} a multi-inquilino}
\label{sec:evolucao}

O projeto nasceu com uma filosofia \textbf{\emph{local-first}}: toda a
computação e persistência ocorriam na máquina do utilizador, com
autenticação \textsc{jwt} \textsc{hs256} própria, base de dados SQLite e
ficheiros guardados localmente. Este desenho cumpria o princípio de que
\emph{nenhum dado saísse da máquina}, e corresponde ao que foi inicialmente
planeado: SQLite como base de dados estruturada e o sistema de ficheiros
para os \emph{media} originais.

A necessidade de disponibilizar o sistema publicamente motivou uma
\textbf{migração para uma arquitetura de nuvem multi-inquilino}:

\begin{itemize}[leftmargin=*]
  \item A \textbf{autenticação} passou de um fluxo HS256 próprio para o
  \textbf{Supabase Auth}, com tokens assinados por chave ECC P-256 e
  validados contra o \emph{endpoint} JWKS público.
  \item O \textbf{modelo de domínio} deixou de ser de dono único
  (\emph{single-owner}) para \textbf{multi-inquilino}: todas as entidades
  referenciam um
  \texttt{user\_id} (\textsc{uuid}) e todas as consultas filtram por esse
  campo.
  \item O \textbf{armazenamento} migrou do disco local para o
  \textbf{Supabase Storage}, passando a referência de cada ficheiro a ser
  uma chave de objeto em vez de um caminho de disco.
\end{itemize}

É importante sublinhar que a migração para a nuvem \emph{não} eliminou o
modo local: o sistema corre de forma \textbf{idêntica nos dois ambientes}
--- tudo o que se faz \emph{online} é possível \emph{localmente}, e
vice-versa --- partilhando o mesmo código, a mesma base de dados e o mesmo
\emph{Storage} (Supabase). Esta paridade é um requisito explícito do
projeto e tem consequências concretas: por exemplo, os vídeos são
\emph{produzidos} em ambiente local (onde não há limite de memória) e
\emph{reproduzidos} no \emph{site} implantado, por escreverem no mesmo
\emph{Storage} (\Cref{sec:impl_m4}). Quanto à IA, o motor de texto por
omissão é o \textbf{Groq} (nuvem, gratuito e rápido), mas a inferência
\textbf{local} via \textbf{Ollama} permanece disponível como opção de
privacidade e de operação \emph{offline}. Esta evolução é um exemplo
concreto de compromisso maduro entre acessibilidade, privacidade e
escalabilidade.

\section{Fluxo de Dados e Execução}
\label{sec:fluxo}

\subsection{Encadeamento do \emph{pipeline}}

Os quatro módulos formam uma cadeia em que a saída de cada um alimenta o
seguinte. A \Cref{tab:fluxo} resume, para cada módulo, o que consome e o
que produz — tornando explícito porque é que a ordem importa (por exemplo,
o M3 só pode ancorar a narrativa em factos depois de o M1 os ter
extraído).

\begin{table}[H]
\centering
\caption{Fluxo de dados ao longo do \emph{pipeline} M1--M4.}
\label{tab:fluxo}
\small
\begin{tabularx}{\linewidth}{@{}p{1.5cm}YY@{}}
\toprule
\textbf{Módulo} & \textbf{Consome} & \textbf{Produz} \\
\midrule
M1 & Bytes de fotografias, documentos e GEDCOM & \texttt{MediaFile} (metadados
\textsc{exif} + descrição de IA + \textsc{ocr}); \texttt{Person} e relações. \\
M2 & \texttt{MediaFile} e \texttt{Person} & Datas resolvidas;
\texttt{TimelineEvent}; grafo de parentesco. \\
M3 & Factos do M1/M2 (via RAG) + intenção do utilizador & \texttt{Story}
(prosa + \texttt{scenes} ligadas a fotografias). \\
M4 & \texttt{Story} (cenas) + fotografias do \emph{Storage} &
\texttt{VideoOutput} (MP4 sincronizado, no \emph{Storage}). \\
\bottomrule
\end{tabularx}
\end{table}

\subsection{Sequência de uma tarefa em segundo plano}

Quando uma operação corre em segundo plano (em produção, a análise visual
diferida das fotografias do M1), o fluxo é o seguinte:

\begin{enumerate}[leftmargin=*]
  \item A API valida o pedido, cria um registo \texttt{TaskRecord} (estado
  \texttt{pending}) e devolve \emph{imediatamente} o seu identificador.
  \item A tarefa é despachada para o Celery ou para o executor
  \emph{in-process} (\Cref{sub:async}), que a marca \texttt{running}.
  \item O corpo da tarefa executa; ao terminar, o registo passa a
  \texttt{done} (ou \texttt{failed}) e é ligado ao artefacto produzido.
  \item O \emph{frontend} faz \emph{polling} do estado e atualiza a
  interface na transição para \texttt{done}.
\end{enumerate}

A geração de \textbf{narrativa} e de \textbf{vídeo} pode, em código, ser
pedida nesta modalidade (\texttt{mode=background}); contudo, na
configuração implantada corre \textbf{sincronamente} (\Cref{sub:async},
\Cref{sec:impl_m4}), mantendo a instância acordada durante o processamento
e devolvendo o resultado diretamente --- evitando que uma instância gratuita
que adormeça abandone a tarefa a meio.

\section{Modelo de Dados}
\label{sec:modelo_dados}

A camada de persistência relacional é modelada com SQLAlchemy e materializada
em PostgreSQL (Supabase). As entidades principais são:

\begin{itemize}[leftmargin=*]
  \item \texttt{MediaFile} — cada ficheiro carregado (fotografia,
  documento, vídeo). Guarda o nome original, a chave de objeto no Storage,
  os metadados \textsc{exif} (data, GPS, câmara), os campos de análise de
  IA (\texttt{ai\_description}, \texttt{ai\_setting}, \texttt{ai\_emotion},
  \texttt{ai\_tags}, \texttt{ai\_narrative\_hint}), o texto de \textsc{ocr},
  o estado de processamento e \texttt{person\_ids} (as pessoas da árvore
  etiquetadas na fotografia).
  \item \texttt{Person} — cada pessoa, importada do GEDCOM \emph{ou criada
  manualmente} (nome, sexo, datas, locais, identificador GEDCOM, etiqueta
  de família).
  \item \texttt{Relationship} — uma aresta de parentesco entre duas
  pessoas (\texttt{pai}/\texttt{mãe}/\texttt{cônjuge}). As relações deixaram
  de viver apenas num grafo JSON em disco (efémero) e passaram a ser
  persistidas na base de dados, sustentando o editor manual e a vista
  interativa da árvore.
  \item \texttt{TimelineEvent} — eventos da linha temporal (título,
  descrição, tipo, local, data resolvida).
  \item \texttt{Story} — narrativa gerada (título, texto, tipo de evento,
  \emph{template} usado, \emph{backend} de \textsc{llm}, número de factos
  usados, \emph{prompt}, idioma, projeto associado). Guarda ainda as
  \emph{preferências de produção} do vídeo --- voz (\texttt{voice}),
  legendas (\texttt{subtitles}, \texttt{subtitle\_size}) ---, a lista de
  fotografias explicitamente escolhidas (\texttt{media\_ids}), um indicador
  de \emph{favorita} (\texttt{favorite}) e o campo \texttt{scenes}
  (\textsc{json}): a forma segmentada em cenas da narrativa --- uma lista de
  \texttt{\{text, photo\_ids, caption\}} que permite ao M4 sincronizar cada
  fotografia com a sua narração (\Cref{sub:cenas}).
  \item \texttt{VideoOutput} — vídeo gerado (nome, chave de objeto,
  tamanho, fotografias usadas, estado, história e projeto associados), bem
  como a referência à \emph{faixa de legendas} \texttt{.vtt} e o tamanho de
  legenda escolhido, consumidos pelo leitor (\Cref{sub:narracao_legendas}).
  \item \texttt{Project} e \texttt{ProjectMedia} — projetos (coleções) e
  a relação muitos-para-muitos com as fotografias. Cada projeto é um
  espaço de trabalho \emph{isolado}: as suas fotografias, a sua linha
  temporal (derivada apenas dessas fotografias), a sua família/árvore
  (identificada pela etiqueta do projeto), as suas histórias e vídeos não
  se misturam com os de outros projetos nem com a Biblioteca global.
  \item \texttt{TaskRecord} — registo de tarefas assíncronas (estado,
  \emph{payload}, carimbos temporais, erro). Cada registo guarda também o
  \texttt{project\_id} a que pertence, de modo que as tarefas em segundo
  plano de um projeto não se misturam com as dos restantes (filtragem por
  projeto na página de Tarefas).
\end{itemize}

As relações principais entre entidades resumem-se na \Cref{tab:relacoes}.
Uma fronteira de tipos relevante: o \texttt{user\_id} é um \textsc{uuid}
(proveniente do Supabase Auth), ao passo que as chaves primárias das
entidades de domínio são inteiros sequenciais do PostgreSQL. A
\Cref{fig:er} (Anexo~\ref{ap:diagramas}) apresenta o diagrama
entidade-relação.

\begin{table}[H]
\centering
\caption{Relações principais entre entidades.}
\label{tab:relacoes}
\small
\begin{tabular}{@{}lcl@{}}
\toprule
\textbf{Entidade A} & \textbf{Cardinalidade} & \textbf{Entidade B} \\
\midrule
Project      & 1 --- N & ProjectMedia \\
ProjectMedia & N --- 1 & MediaFile \\
Story        & 1 --- N & VideoOutput \\
Project      & 1 --- N & Story \\
TaskRecord   & N --- 1 & Story / VideoOutput (resultado) \\
\bottomrule
\end{tabular}
\end{table}


\section{Design de Interface (UI/UX)}
\label{sec:uiux}

\subsection{Princípios e sistema de design}

O \emph{design} visual procura evocar o \textbf{calor do arquivo
familiar}. Assenta numa paleta \emph{amber}/\emph{stone} (tons quentes
que remetem para fotografias antigas e luz de candeeiro), com acentos
secundários sóbrios. A tipografia é dupla: \emph{Inter} para a interface
(legibilidade) e \emph{Fraunces} (serifada) para títulos e para a prosa
do leitor de histórias, conferindo um registo literário. As animações são
discretas (\emph{fade-in}, \emph{slide-up}, \emph{skeletons} de
carregamento), realizadas com animações \textsc{css} declaradas no Tailwind.

Um detalhe de implementação relevante para a experiência é o
\textbf{tema} (claro/escuro/sistema): um \emph{script} embutido no
cabeçalho lê a preferência guardada \emph{antes} do primeiro
\emph{paint}, evitando o desagradável ``flash'' de tema errado ao
carregar a página.

\subsection{Estrutura de navegação}

A aplicação organiza-se em rotas \textbf{públicas} (\emph{landing}/login,
registo, recuperação de palavra-passe, ``sobre'') e \textbf{protegidas},
acessíveis após autenticação através de uma barra lateral
(\emph{sidebar}) responsiva. As secções visíveis são: \emph{Início},
\emph{Biblioteca}, \emph{Família}, \emph{Linha temporal}, \emph{Projetos},
\emph{Gerar}, \emph{Histórias}, \emph{Vídeos} e \emph{Definições}. A página
de \emph{Tarefas} permanece acessível pela rota \texttt{/tasks}, mas foi
\textbf{retirada da barra lateral} por as gerações passarem a correr de
forma síncrona, tornando o acompanhamento por \emph{polling} dispensável no
uso corrente.

As capturas de ecrã dos ecrãs centrais da aplicação --- \emph{Dashboard},
Biblioteca, assistente \emph{Gerar}, leitor de histórias e galeria de
vídeos --- constam do Anexo~\ref{ap:capturas}.

\subsection{Internacionalização e acessibilidade}

Toda a interface está disponível em \textbf{português europeu} (omissão)
e \textbf{inglês}, com a preferência persistida e o atributo
\texttt{<html lang>} sincronizado. Ao nível da acessibilidade, existem
rótulos \textsc{aria} e navegação por teclado (por exemplo, no
\emph{lightbox} da biblioteca); uma auditoria formal de conformidade
\textsc{wcag} é apontada como trabalho futuro (\Cref{cap:conclusoes}).

% =====================================================================
%  CAPÍTULO 6 — IMPLEMENTAÇÃO
% =====================================================================
\chapter{Implementação}
\label{cap:implementacao}

Este capítulo descreve como os requisitos e o \emph{design} foram
traduzidos em código. Organiza-se pelos quatro módulos do \emph{pipeline}
(M1--M4), seguidos da infraestrutura do \emph{backend}, do \emph{frontend}
e das decisões de segurança. Ao longo do texto incluem-se excertos de
código \emph{críticos} — não a totalidade da implementação, mas os pontos
que melhor ilustram as decisões técnicas.

\section{Módulo M1 --- Ingestão Multimodal}
\label{sec:impl_m1}

O M1 é o ponto de entrada de todo o material. Recebe os \emph{bytes} de um
ficheiro e orquestra uma cadeia de processamento: \emph{security scan}
$\rightarrow$ \textsc{exif} $\rightarrow$ carregamento para o \emph{Storage}
$\rightarrow$ análise visual $\rightarrow$ \textsc{ocr}, persistindo um
registo \texttt{MediaFile}.

O processamento decorre sobre um ficheiro temporário (as bibliotecas de
\textsc{exif}/\textsc{ocr}/visão esperam um caminho de sistema), removido no
final. Por segurança, os \emph{bytes} só chegam ao \emph{Storage} \emph{após}
um \emph{scan} bem-sucedido. As etapas são:

\begin{description}[leftmargin=*]
  \item[\emph{Security scan}.] Cálculo de \emph{checksum} MD5 e \emph{scan}
  antivírus com o \textbf{ClamAV} (quando presente no sistema; na sua
  ausência, o passo degrada graciosamente e regista um aviso). Perante a
  deteção de uma ameaça, o registo é marcado \texttt{FAILED} e o
  \emph{pipeline} interrompe-se --- o conteúdo nunca chega ao \emph{Storage}.
  \item[\textsc{exif}.] Data de captura, coordenadas GPS, marca e modelo
  da câmara, extraídos com a biblioteca \textbf{exifread}~\cite{exifread} (puro Python, sem
  dependências de sistema, funcionando de forma idêntica em ambiente local
  e na nuvem) e preservando o \textsc{exif} em bruto. As coordenadas, no
  formato \emph{graus/minutos/segundos}, são convertidas para decimal com o
  sinal correto segundo a referência (N/S, E/W).
  \item[Análise visual.] O \textbf{Gemini Vision} produz descrição da
  cena, contagem de pessoas, cenário, emoção dominante, \emph{tags} e uma
  \emph{sugestão narrativa} — a matéria-prima factual do M3.
  \item[\textsc{ocr}.] Para documentos, o \textbf{Tesseract}~\cite{smith2007tesseract} extrai
  texto, enriquecendo o contexto factual.
\end{description}

\textbf{Análise diferida.} A análise visual é o passo mais lento (uma
chamada ao Gemini, da ordem dos segundos). Para não bloquear o carregamento
--- sobretudo em lotes ---, o M1 corre no pedido HTTP apenas a parte rápida
(\emph{scan}, \textsc{exif} e \emph{Storage}) e \emph{difere} a IA
(Gemini/\textsc{ocr}) para segundo plano, no executor \emph{in-process}
(\Cref{sec:impl_backend}). A fotografia fica logo visível; o registo
permanece em \texttt{PROCESSING} e a descrição preenche-se no fim,
sinalizada por um indicador ``a analisar'' na interface, que faz
\emph{polling} enquanto houver registos por concluir.

\textbf{Re-análise sob procura.} Como o executor de segundo plano pode ser
terminado pela plataforma gratuita antes de concluir (deixando fotografias
presas em \texttt{PROCESSING}), um \emph{endpoint} de \textbf{re-análise}
--- e um botão na Biblioteca --- reprocessa de forma síncrona as que ficaram
sem descrição, marcando-as \texttt{COMPLETED} e desbloqueando a linha
temporal e o vídeo. É \emph{consciente da quota} (indica quantas serão
analisadas antes de avançar) e processa em \textbf{pequenos lotes} por
pedido, para não esgotar a memória da instância gratuita.

A árvore genealógica é tratada por um \emph{parser} \textbf{GEDCOM}~\cite{gedcom}
dedicado (o maior módulo, $\approx$430 linhas), que converte o ficheiro em
registos \texttt{Person} e nas respetivas relações, com suporte a
etiquetagem de famílias (permitindo distinguir várias árvores na mesma
conta).

\section{Módulo M2 --- Organização Temporal}
\label{sec:impl_m2}

O M2 transforma factos dispersos numa estrutura temporal e relacional
coerente, através de dois componentes.

O \textbf{\emph{date resolver}} normaliza datas heterogéneas e
frequentemente incompletas, produzindo uma chave de ordenação consistente
— essencial, dado que o material de arquivo raramente tem datas exatas.
Trata os qualificadores do formato GEDCOM (\texttt{ABT} ``cerca de'',
\texttt{BEF} ``antes de'', \texttt{AFT} ``depois de'') e os casos de apenas
mês/ano ou só ano, e rejeita datas implausíveis (anteriores a 1800, o
\emph{epoch} Unix de 1970 ou no futuro), evitando que ruído contamine a
linha temporal. Os casamentos importados do GEDCOM geram automaticamente
\texttt{TimelineEvent}s datados, que passam a fazer parte do contexto
factual disponível ao M3.

O coração do M2 é o \textbf{grafo de parentesco} (\texttt{FamilyGraph}),
construído sobre o \emph{NetworkX}~\cite{networkx} como grafo dirigido: os
nós são pessoas e as arestas representam relações (\emph{parent},
\emph{child}, \emph{spouse}, \emph{sibling}). A classe produz
\emph{resumos narrativos} do parentesco (por exemplo, \emph{``João (pai de
Maria, marido de Ana)''}), consumidos diretamente pelo M3 para dar ao
\textsc{llm} o contexto de \emph{quem é quem}.

\textbf{Relações na base de dados e edição manual.} As relações passaram a
ser persistidas numa tabela \texttt{relationships} --- a fonte de verdade
durável ---, em vez de um grafo JSON em disco que, sendo efémero na nuvem,
deixava o contexto familiar vazio após cada reinício da instância (uma
fonte subtil de narrativas sem relações). Isto sustenta, na página
\emph{Família}, uma \textbf{vista interativa da árvore} (\emph{React~Flow},
\emph{layout} por gerações, casais lado a lado, nós coloridos por sexo e
ligações desenhadas com um \textbf{barramento} partilhado --- descida do
casal, barra de irmãos e ligação curta a cada filho ---, evitando linhas
sobrepostas) e um \textbf{editor manual} de pessoas e ligações
(\emph{pai}/\emph{mãe}/\emph{filho}/\emph{irmão}/\emph{casado(a)}), com um
assistente de \textbf{criação rápida de árvore} (eu, pais, avós, irmãos).
Após cada edição, o \texttt{FamilyGraph} é reconstruído da base de dados,
isolado \emph{por utilizador}, mantendo o M3 sincronizado.

As \textbf{notas} de cada pessoa entram no contexto do M3, influenciando a
narrativa. Cada fotografia pode ser \textbf{etiquetada com quem nela
aparece} (\texttt{media\_files.person\_ids}), ligando \emph{rostos a nomes}
--- a ponte entre a árvore e a narrativa ---, e cada pessoa pode ter uma
\textbf{foto de perfil} (\texttt{persons.photo\_media\_id}), usada como
\emph{avatar} na lista e nos nós da árvore. A árvore pode ainda ser
\textbf{exportada para GEDCOM~5.5.1} (re-importável noutras plataformas),
demonstrando \emph{interoperabilidade} com o formato standard.

\textbf{Isolamento de múltiplas árvores.} Importar vários ficheiros GEDCOM
levantava um problema subtil: as exportações genealógicas reiniciam quase
sempre os identificadores em \texttt{@I1@} por ficheiro, pelo que a
deduplicação por \texttt{(user\_id, gedcom\_id)} fazia a segunda família
\emph{sobrescrever} a primeira, fundindo árvores sem relação. A solução foi
\textbf{prefixar o \texttt{gedcom\_id} persistido com a etiqueta da família}
(\texttt{"<família>::<id>"}; sem etiqueta, um \emph{id} de lote único), de
modo que reimportar a \emph{mesma} família continua idempotente mas famílias
distintas nunca colidem. Sobre este mecanismo, cada \textbf{projeto} mantém
a sua família isolada e, dentro dele, cada GEDCOM importado fica num
\textbf{sub-grupo próprio} (nomeado pelo ficheiro), alternável por
\emph{chips} — uma vista por sub-família ou todas em conjunto (parâmetro
\texttt{group} na \textsc{api}, que casa a etiqueta do projeto e as suas
sub-etiquetas).

\section{Módulo M3 --- Geração Narrativa}
\label{sec:impl_m3}

O M3 é o módulo de IA generativa central. Combina três componentes: um
\emph{cliente LLM} híbrido, o subsistema \emph{RAG} e um conjunto de
\emph{templates}.

\subsection{Cliente LLM híbrido e \emph{fallback}}

O \texttt{LLMClient} implementa uma \textbf{cadeia de \emph{fallback} de
três níveis} para o texto: o motor primário é o \emph{Groq}
(\emph{Llama~3.3~70B}, camada gratuita rápida); na sua ausência ---
tipicamente sem Internet --- recorre-se ao \emph{Ollama} local
(Llama~3.2~3B), \emph{offline e privado}; e só se ambos falharem entra o
\emph{Gemini~Flash}. A ordem é \textbf{deliberada}: o Llama~3.2~3B local, em
\textsc{cpu}, é demasiado lento para uso interativo, pelo que o Groq (70B na
nuvem) dá melhor qualidade e resposta quase imediata; e a quota do Gemini
fica \textbf{reservada para a visão} (M1), que se mantém \emph{Gemini-only}
por o Groq não oferecer visão --- repartindo texto e visão por dois
serviços gratuitos sem competirem pela mesma quota. Como mostra o
\Cref{lst:llm}, se um nível falhar \emph{durante} a geração há
\emph{fallback} em tempo de execução; só quando \emph{todos} falham é
lançada uma exceção que marca a tarefa como falhada, evitando guardar uma
narrativa falsa.

\subsection{Retrieval-Augmented Generation com degradação graciosa}

O \texttt{RAGSystem} indexa os factos extraídos pelo M1/M2 numa base
vetorial (ChromaDB~\cite{chromadb}), com o \texttt{user\_id} na
\emph{metadata} de cada documento; a pesquisa filtra sempre por esse campo
(isolamento multi-inquilino). Quando a biblioteca é grande e o utilizador
indica um foco, o método \texttt{search\_media\_ids} injeta no \emph{prompt}
apenas o subconjunto de fotografias mais relevante ao tema (por
similaridade vetorial), mantendo o contexto focado, dentro da janela do
modelo e menos sujeito a alucinação.

A decisão de engenharia mais interessante é a \textbf{degradação graciosa}.
Quando o ChromaDB não está instalado --- por exemplo, no \emph{deploy}
gratuito do Render, onde o custo de \emph{build} é incomportável ---, uma
importação condicional deteta a ausência e comuta o \texttt{RAGSystem} para
um \emph{stub}: a indexação torna-se inócua e a pesquisa devolve listas
vazias, passando o M3 a usar a totalidade dos factos. O sistema mantém-se
\emph{funcional}, perdendo apenas o passo de \emph{retrieval}.

\subsection{\emph{Templates} e calibração para português europeu}

Existem \textbf{sete} \emph{templates}: seis temáticos (\emph{Memória
Familiar}, \emph{Momento Capturado}, \emph{História de Amor},
\emph{Aventura Familiar}, \emph{Nova Vida}, \emph{Momento de Alegria}) e
um \emph{livre} (\emph{personalizado}), em que o utilizador define tom e
estrutura. Cada \emph{template} combina um \emph{prompt} com metadados de
apresentação.

O desafio linguístico central é que os modelos disponíveis tendem para o
\textbf{português do Brasil}. A resposta foi a constante
\texttt{PT\_PT\_RULES}, embutida em todos os \emph{templates}, com regras
explícitas de léxico, colocação pronominal e proibição de gerúndio
(\Cref{lst:ptpt}). Cada \emph{template} recebe ainda um bloco
\texttt{USER\_FOCUS\_BLOCK} de \emph{prioridade máxima}: a intenção
textual do utilizador determina o \emph{assunto}, enquanto os factos
servem de matéria-prima. O idioma de saída pode ser comutado para inglês
através de um sufixo de \emph{prompt}, sem \emph{templates} paralelos.

\textbf{Rede de segurança determinística.} Como a instrução por si só não
é fiável --- mesmo avisados, o Llama e o Gemini reincidem no português do
Brasil ---, acrescentou-se um \textbf{pós-processador determinístico}
(\texttt{pt\_pt.py}) que corre \emph{depois} da geração e corrige
mecanicamente os brasileirismos mais frequentes: substituições de
vocabulário inequívocas e que preservam género/número (\emph{trem}\,$\rightarrow$\,\emph{comboio},
\emph{fazenda}\,$\rightarrow$\,\emph{quinta}, \emph{caminhão}\,$\rightarrow$\,\emph{camião}\ldots)
e a conversão da construção \emph{auxiliar~+~gerúndio} para
\emph{auxiliar~+~a~+~infinitivo} (``estava chorando''\,$\rightarrow$\,``estava a
chorar''), o sinal gramatical mais reconhecível da variante. O mesmo
módulo expõe uma função \texttt{count\_brasileirismos} usada como métrica
objetiva de aderência ao PT-PT (\Cref{sec:qualidade}). O pós-processador
corrige ainda a \textbf{concordância de género} num conjunto curado de
substantivos (``o~lã''\,$\rightarrow$\,``a~lã''), \textbf{remove parágrafos
repetidos} (que o modelo por vezes produz para encher extensões longas) e
\textbf{corta uma frase final incompleta} deixada por um limite de
\emph{tokens}.

\textbf{Estilo e integração das imagens.} Dois blocos de \emph{prompt}
elevam a qualidade literária: as \texttt{STYLE\_RULES} proíbem a repetição
de nomes próprios (impondo pronomes e perífrases), exigem variedade de
verbos e conectores e reforçam a colocação pronominal europeia; as
\texttt{MEDIA\_GROUNDING} instruem o modelo a situar o espaço, o lugar e as
pessoas \emph{antes} de descrever a cena, com transições suaves, para que
cada fotografia surja em contexto e não de forma abrupta.

\textbf{Controlo de extensão.} O assistente permite escolher a duração em
quatro níveis --- \emph{curta} ($\approx$1\,min), \emph{média}
($\approx$2--3), \emph{longa} ($\approx$4--5) e \emph{épica}
($\approx$6--8). Cada nível gera uma instrução de extensão prioritária (que
sobrepõe o número de parágrafos do \emph{template}) e um limite de
\emph{tokens} calibrado para a duração ($\approx$140~palavras/minuto). Para
não sacrificar a qualidade, a \textbf{não-repetição tem prioridade sobre o
comprimento}: se os factos não bastarem, o modelo escreve um texto mais
curto mas completo.

\subsection{Segmentação em cenas (sincronização narrativa--imagem)}
\label{sub:cenas}

Para que o vídeo deixe de ser um \emph{slideshow} com tempos arbitrários
e passe a um verdadeiro documentário, introduziu-se uma etapa de
\textbf{segmentação em cenas} (\texttt{scene\_builder.py}). Após gerar a
prosa, o M3 divide-a em parágrafos (cada parágrafo $=$ uma cena) e
associa a cada cena as fotografias que melhor a ilustram, produzindo uma
estrutura serializável guardada na coluna \texttt{stories.scenes} --- uma
lista de \texttt{\{text, photo\_ids, caption\}}, cuja forma se ilustra no
\Cref{lst:scene} (Anexo~\ref{ap:codigo}).

A atribuição foto\,$\rightarrow$\,parágrafo é \emph{determinística e sem
dependências externas}: pontua-se a sobreposição léxica entre cada
parágrafo e a ``impressão digital'' textual de cada fotografia (descrição
de IA, \emph{tags}, cenário, sugestão narrativa, \textsc{ocr}, local). É,
na prática, um passo de \emph{retrieval} que ancora os visuais à prosa.
Quando não há sinal léxico, distribuem-se cronologicamente. Numa narrativa
\textbf{longa ilustrada por poucas fotografias}, os parágrafos que sobram
--- onde a narração já não se refere a uma imagem em concreto --- são
preenchidos \textbf{alternando} ciclicamente as disponíveis (sem repetir a
da cena anterior), de modo que o vídeo continua a mudar de imagem em vez de
congelar; cada fotografia surge primeiro no parágrafo que a descreve e só
depois é reaproveitada. Toda a etapa é
defensiva: qualquer falha deixa \texttt{scenes\,=\,None} e o M4 regressa ao
\emph{slideshow} clássico. Esta abordagem responde diretamente ao que o
estudo inicial identificou como o desafio técnico central do M4
(\Cref{sub:geracao_video}): cada fotografia visível exatamente durante a
sua narração.

\subsection{Ancoragem nas fotografias e seleção pelo utilizador}
\label{sub:selecao_fotos}

A narrativa \textbf{não} se baseia apenas na árvore e nas notas que o
utilizador escreve: o construtor de contexto (\texttt{\_build\_events\_from\_media})
incorpora, para cada fotografia, a \textbf{descrição gerada pela visão}, o
\textsc{ocr} (texto de documentos, como certidões), o cenário (o
\emph{espaço} da cena) e o \textbf{lugar} (topónimo derivado do GPS), a
emoção, as \emph{tags} e \emph{quem aparece} na imagem. São estes factos visuais que
ancoram o texto a momentos concretos --- e, por via da segmentação em cenas
descrita acima, determinam também a \textbf{ordem das imagens no vídeo}.

Por defeito, são consideradas todas as fotografias disponíveis (ou, numa
biblioteca grande, as que o \textsc{rag} julga relevantes ao tema). Foi
acrescentada a possibilidade de o utilizador \textbf{escolher explicitamente
quais as fotografias} que entram numa história, diretamente no assistente de
geração: a lista de \texttt{media\_ids} selecionados é enviada ao gerador,
que restringe o contexto a essas imagens (saltando o \emph{retrieval}
automático). Como a narrativa alimenta a segmentação em cenas, esta escolha
propaga-se ao vídeo. O mecanismo funciona tanto no arquivo geral como
\emph{dentro de um projeto} (onde a seleção se limita às fotografias do
projeto), dando ao utilizador controlo curatorial sobre o produto final. Um
projeto pode ainda \textbf{reutilizar fotografias já analisadas} da
biblioteca geral, ligando-as sem novo carregamento nem nova análise --- a
análise de visão é feita uma única vez e reaproveitada ---, para além de
poder receber fotografias carregadas diretamente para si.

A seleção é ainda \textbf{organizada por pessoa}: ao escolher os
intervenientes da história, as fotografias surgem agrupadas por cada pessoa
(as imagens em que está etiquetada, via \texttt{media\_files.person\_ids}),
permitindo escolher exatamente quais das fotos de cada um entram. Em
complemento, cada pessoa tem uma \textbf{galeria própria} onde se anexam ou
removem as suas fotografias e documentos (por exemplo, o retrato e uma
certidão), gerindo o conjunto do lado da pessoa --- o inverso da etiquetagem
``quem aparece'' feita do lado da fotografia.

\subsection{Regeneração orientada por \emph{feedback}}
\label{sub:regenerar}

A geração raramente acerta à primeira no registo pretendido. Para o permitir sem obrigar a recomeçar, o leitor
de histórias oferece uma
\textbf{regeneração \emph{in-place} orientada por \emph{feedback}}: o
utilizador escreve uma nota em texto livre (por exemplo, ``mais curta'',
``foca-te no avô João'' ou ``tom mais sóbrio'') e o sistema reescreve a
narrativa \emph{sobre a mesma história}, reutilizando as mesmas
fotografias, pessoas, idioma e definições de voz/legendas, e preservando o
mesmo identificador (a nota é incorporada no \emph{prompt} como instrução
de prioridade máxima). O mecanismo reusa o gerador do M3 com um parâmetro
\texttt{update\_story\_id}, atualizando o texto e a respetiva segmentação
em cenas em vez de criar um registo novo --- o que evita acumular versões
e mantém o vídeo eventualmente já gerado coerente com a última narrativa.

\section{Módulo M4 --- Geração Multimédia}
\label{sec:impl_m4}

O M4 converte uma narrativa num \textbf{vídeo documental} MP4 (H.264,
$1280\times720$, 24~fps por defeito --- a resolução e os \emph{frames}/segundo
são parametrizáveis por variável de ambiente). O \texttt{M4Processor} orquestra
a seleção cronológica das fotografias, a síntese de voz da narrativa e a
montagem do vídeo, mantendo apenas espaço de trabalho temporário no disco (as
fotografias e o vídeo final vivem no \emph{Storage}). A montagem corre numa
\emph{thread} dedicada e o pedido é tratado de forma \textbf{síncrona}, mantendo
a instância acordada durante o \emph{render}.

\textbf{Geração local, reprodução universal.} A montagem com MoviePy/FFmpeg
é intensiva em memória e a 720p excede os $\approx$512~MB do plano gratuito do
Render (\emph{out-of-memory}, que deixava o vídeo preso em processamento). Em
vez de degradar a qualidade, optou-se por uma política de \textbf{geração
exclusivamente local}: na nuvem, a flag \texttt{VIDEO\_LOCAL\_ONLY} faz o
\emph{backend} \textbf{recusar} de imediato o pedido de \emph{render} (com uma
mensagem clara), ao passo que, em ambiente local, o vídeo é montado a 720p sem
restrição de memória. Como o \emph{backend} local escreve o MP4 no \emph{mesmo}
\emph{Storage} (Supabase) que a nuvem consome, o vídeo \textbf{pré-gerado
localmente} fica disponível e reproduzível no site implantado sem qualquer
\emph{render} em tempo real --- a reprodução é leve (apenas servir o ficheiro) e
nunca é afetada pelo limite de memória. Esta separação entre \emph{produzir}
(local, pesado) e \emph{reproduzir} (nuvem, leve) é uma consequência direta do
requisito de que tudo o que se faz \emph{online} seja também possível
\emph{localmente}, e vice-versa.

\subsection{Montagem cinematográfica}

O \texttt{video\_builder} eleva a qualidade percebida acima de um simples
\emph{slideshow}:

\begin{itemize}[leftmargin=*]
  \item \textbf{Enquadramento não-destrutivo} (\emph{letterbox}): cada
  fotografia cabe inteira no quadro, sem cortes nem distorção, sobre um
  fundo composto pela mesma imagem ampliada, desfocada (\emph{Gaussian
  blur}) e escurecida — um pano de fundo cinematográfico.
  \item \textbf{Movimento por fotografia configurável}
  (\texttt{VIDEO\_MOTION}): \emph{none} (estático), \emph{gentle} ou o
  \textbf{efeito Ken~Burns} --- \emph{zoom} subtil ($1.0\!\rightarrow\!1.06$)
  e \emph{pan} direcional com aceleração suave (\emph{smoothstep}), com o
  ângulo determinístico a partir do \emph{hash} do caminho do ficheiro para
  que fotografias consecutivas não se movam na mesma direção
  (\Cref{lst:kenburns}). Por \textbf{omissão} usa-se o modo \emph{none}
  (fotografias estáticas): além de aliviar a montagem, permite ``pré-cozer''
  uma única vez cada fotograma de cena, acelerando o \emph{render}.
  \item \textbf{Transições} por \emph{crossfade} e \emph{fade-to-black}
  final, evitando cortes abruptos. O \emph{crossfade} usa uma implementação
  própria (\texttt{\_concatenate\_with\_crossfade}) que mistura apenas os
  \emph{frames} sobrepostos, em vez do \texttt{CompositeVideoClip} do
  MoviePy --- uma otimização que reduziu drasticamente o tempo de montagem
  (\Cref{sec:resultados_testes}).
  \item \textbf{Cartões de título e de fim} de abertura e fecho. O restante
  vídeo é mantido \textbf{limpo}, sem texto gravado na imagem: as legendas
  de narração deixaram de ser ``cozidas'' como \emph{lower thirds} e passaram
  a uma \textbf{faixa de legendas comutável} (ver
  \Cref{sub:narracao_legendas}), mais legível e opcional.
  \item \textbf{Ritmo calibrado}: piso de $\sim$4~s por fotografia em modo
  cena (subido de 3~s), para que nenhuma imagem passe a correr contra o
  \emph{crossfade} de 1.6~s.
\end{itemize}

\subsection{Narração, voz e legendas}
\label{sub:narracao_legendas}

A narração é gerada por \textbf{síntese de voz} (vozes neuronais
europeias do \emph{edge-tts}, com \emph{gTTS} como alternativa) na língua
em que a história foi escrita (capturada em \texttt{Story.language}),
garantindo coerência mesmo que o utilizador altere depois o idioma da
interface. O utilizador pode \textbf{escolher a voz} --- masculina ou
feminina --- no assistente de geração (\texttt{Story.voice}); o M4 resolve
essa preferência para a voz neuronal correspondente em cada idioma.

Quando a história tem cenas (\Cref{sub:cenas}), o M4 entra em \textbf{modo
documentário} (\texttt{build\_documentary}): sintetiza \emph{uma narração
por cena} e mostra as fotografias dessa cena durante (aproximadamente) a
duração da sua narração, concatenando os áudios em ordem após o cartão de
título. Assim, o espetador vê a fotografia que está a ser narrada --- a
sincronização semântica que distingue um documentário de uma apresentação.
O \Cref{lst:plandur} mostra o cálculo determinístico (e testado) da duração
de cada fotografia. Na ausência de cenas, usa-se o alinhamento clássico por
duração total. A \Cref{fig:frame_video} (Anexo~\ref{ap:capturas}) mostra um
fotograma do resultado.



\textbf{Legendas como faixa sincronizada (WebVTT).} Em vez de gravar o
texto na imagem, a narração é exportada como uma \textbf{faixa WebVTT}
(\texttt{.vtt}) \emph{comutável} no leitor, sincronizada \textbf{ao nível
da palavra}: o \emph{edge-tts} devolve, com o áudio, os carimbos temporais
de cada palavra (\emph{WordBoundary}), a partir dos quais se escreve o
\texttt{.vtt}. É servido por um \emph{endpoint} próprio e, no leitor, uma
\texttt{<track>} liga/desliga as legendas e escolhe o tamanho (via
\texttt{::cue}). Para o descarregamento, são ainda \emph{embutidas} no MP4
(\emph{mux} \texttt{mov\_text} via FFmpeg). Assim, o vídeo fica limpo por
omissão, mas a acessibilidade textual está sempre disponível.

\section{Infraestrutura do \emph{Backend}}
\label{sec:impl_backend}

\subsection{API e organização}

O \emph{backend} expõe uma API REST versionada (\texttt{/api/v1/})
organizada em grupos de \emph{routers}: \texttt{auth}, \texttt{upload},
\texttt{timeline}, \texttt{narrative}, \texttt{projects},
\texttt{genealogy}, \texttt{multimedia}, \texttt{tasks} e \texttt{health}.
A documentação interativa (\emph{Swagger UI}) é gerada automaticamente.

\subsection{Fila de tarefas e \emph{sweep} de órfãs}

Gerações longas podem correr em segundo plano, com o \emph{frontend} a
acompanhar o estado por \emph{polling} de \texttt{/api/v1/tasks/\{id\}}.
A execução é feita por uma de duas vias, transparentes para o cliente:

\begin{itemize}[leftmargin=*]
  \item \textbf{Celery + Redis} quando existe um \emph{worker} dedicado
  (ambiente completo).
  \item \textbf{Executor \emph{in-process}} quando \texttt{CELERY\_ENABLED}
  é falso (por exemplo, no plano gratuito da nuvem, que não comporta um
  \emph{worker} pago). Neste modo, a tarefa corre num pequeno
  \emph{thread pool} limitado dentro do próprio processo do \emph{backend}
  --- cada \emph{thread} com o seu próprio \emph{event loop}, à imagem de um
  \emph{worker} Celery, para não bloquear o \emph{loop} principal do
  FastAPI. O pedido devolve igualmente um \texttt{task\_id} de imediato.
\end{itemize}

Ambas as vias partilham os mesmos \emph{corpos} de tarefa
(\texttt{tasks/bodies.py}) e o mesmo mecanismo de transição de estado
(\texttt{running}\,$\rightarrow$\,\texttt{done}/\texttt{failed}), pelo que a
experiência do utilizador é idêntica em qualquer ambiente. O \emph{thread
pool} in-process limita-se a um trabalhador --- prudente face aos
$\approx$512~MB do plano gratuito.

Uma decisão de robustez relevante é o \textbf{\emph{sweep} de tarefas
órfãs} no arranque: como nem um \emph{worker} Celery nem um \emph{thread}
\emph{in-process} sobrevivem a um reinício do \emph{backend}, qualquer
tarefa que ficasse \texttt{pending}/\texttt{running} de um arranque
anterior é, por definição, órfã, sendo marcada como falhada --- o que evita
o sintoma de tarefas eternamente ``em execução''.

\subsection{Observabilidade e validação}

O sistema usa \textbf{structlog} para \emph{logging} estruturado (stdout +
ficheiro rotativo). A sonda profunda \texttt{/healthz} verifica todas as
dependências (base de dados, Ollama, Groq, Gemini, Redis, ChromaDB, disco). Os
\emph{uploads} são validados por \emph{magic bytes} (libmagic), tipo MIME
e limites de tamanho (25~MB para fotografias, 100~MB para GEDCOM), e o
\emph{rate limiting} por IP (slowapi) protege contra abuso, com limites
diferenciados (\Cref{tab:rate}).

\begin{table}[H]
\centering
\caption{Limites de taxa aplicados pelo \emph{middleware} (por IP).}
\label{tab:rate}
\small
\begin{tabular}{@{}lc@{}}
\toprule
\textbf{Classe de \emph{endpoint}} & \textbf{Limite} \\
\midrule
Por omissão & 100/minuto \\
Carregamento (\emph{upload}) & 20/minuto \\
Geração (narrativa + vídeo) & 5/minuto \\
\bottomrule
\end{tabular}
\end{table}

\section{\emph{Frontend}}
\label{sec:impl_frontend}

A \textsc{spa} React/TypeScript distingue rigorosamente \emph{estado de
servidor} (gerido pelo \textbf{TanStack Query}, com $\approx$50 \emph{hooks}
que encapsulam pedidos, \emph{cache}, invalidação e \emph{polling}) de
\emph{estado de cliente} (gerido pelo \textbf{Zustand}: autenticação e
tema). Um interceptor do cliente HTTP (axios) injeta automaticamente o
\emph{token} \emph{bearer} em cada pedido e, perante uma resposta
\texttt{401}, termina a sessão de forma limpa. Um \texttt{SessionLoader}
hidrata a sessão Supabase no arranque e instala um \emph{listener} que
mantém o \emph{token} fresco, incluindo renovações silenciosas.

A \textbf{navegação} cobre os nove destinos visíveis da aplicação (Início,
Biblioteca, Família, Linha Temporal, Projetos, Gerar, Histórias, Vídeos e
Definições; a página de \emph{Tarefas} mantém-se acessível pela rota, mas
fora da barra lateral), e cada história exibe um \emph{badge} com o
\textbf{projeto} a que pertence (com ligação direta), clarificando a
organização do arquivo. A
\textbf{página pública} (\emph{landing}) inclui uma \textbf{demonstração
interativa} do \emph{pipeline} (foto $\rightarrow$ análise por IA $\rightarrow$
narrativa $\rightarrow$ vídeo), com a descrição e a história a serem
``escritas'' ao vivo, comunicando o valor do sistema antes do registo.

\textbf{Gestão de histórias.} A lista ganhou ferramentas de curadoria:
\textbf{pesquisa} por título/texto, \textbf{filtro} por projeto,
\textbf{ordenação} e \textbf{favoritos} (estrela que fixa as preferidas no
topo), além de metadados por cartão (fotografias e pessoas usadas). No
leitor, cada história pode ser \textbf{exportada para PDF} e
\textbf{regenerada com \emph{feedback}} (\Cref{sub:regenerar}).

\textbf{Árvore familiar interativa.} A vista \emph{React~Flow}
(\Cref{sec:impl_m2}) permite \textbf{pesquisar e centrar} numa pessoa,
\textbf{destacar a linhagem} ao clicar num nó (o resto da árvore
esbate-se), \textbf{exportar} a árvore como PNG e \textbf{adicionar
familiares} num só passo. Toda a interface está disponível em
\textbf{português europeu} e \textbf{inglês} (\Cref{sec:uiux}), incluindo as
cadeias geradas no cliente.

\subsection{Resiliência a \emph{cold starts}}
\label{sub:coldstart}

O plano gratuito do Render suspende o \emph{backend} ao fim de alguns
minutos de inatividade; o primeiro pedido após a suspensão acorda-o, mas
pode demorar o suficiente para a ligação cair antes de a resposta chegar
ao \emph{browser} — mesmo tendo a operação sido concluída e persistida no
servidor. O \emph{frontend} distingue este caso (resposta perdida, sem
código HTTP) de um erro real e, em vez de mostrar um ``Network Error''
bruto, \emph{reverifica}: aguarda alguns segundos, refaz a consulta
relevante e, se os dados apareceram (uma nova história, novas pessoas
importadas), trata a operação como bem-sucedida. Nunca reenvia o pedido,
para não duplicar dados. Como mitigação de raiz, um \emph{ping} periódico
externo ao \texttt{/healthz} mantém o serviço acordado.
\label{sec:impl_seguranca}

\subsection{Autenticação por JWKS e isolamento multi-inquilino}

A autenticação é delegada ao \textbf{Supabase Auth}. Os tokens são
validados no \emph{backend} contra o \emph{endpoint} JWKS (algoritmo
ES256, chave ECC P-256), com \emph{cache} do documento de chaves em
memória (TTL de 10~min, indexada por \texttt{kid} para suportar rotação de
chaves) e verificação da \emph{audience} (\Cref{lst:jwks}). A identidade
extraída expõe o \textsc{uuid} do utilizador, que é a chave de isolamento:
\textbf{todas} as consultas de domínio filtram por \texttt{user\_id}.

\textbf{Eliminação de conta.} A remoção definitiva de uma conta é feita por um
\emph{endpoint} dedicado que apaga todas as linhas do dono (média, pessoas,
relações, histórias, vídeos, tarefas e projetos) e, em seguida, remove o
utilizador através da \emph{Admin~API} do Supabase (chave \emph{service role}),
garantindo que as credenciais deixam efetivamente de funcionar --- ao contrário
de um mero \emph{sign-out}, que deixava a conta acessível.

\subsection{Servir ficheiros com autenticação}

Os \emph{endpoints} que servem fotografias e vídeos aceitam o \textsc{jwt}
no cabeçalho \texttt{Authorization} \emph{ou} num parâmetro
\texttt{?token=} — necessário porque as \emph{tags} \texttt{<img>} e
\texttt{<video>} não permitem cabeçalhos personalizados. O \emph{backend}
responde com um \emph{redirect} para um URL assinado e temporário do
Supabase Storage, evitando expor as chaves de objeto.

\subsection{Privacidade por inferência local}

A herança \emph{local-first} mantém-se na \textbf{opção} de executar o
\textsc{llm} localmente (Ollama): com o Ollama disponível e sem ligação à
nuvem, a cadeia recai sobre ele e os factos do arquivo não são enviados a
terceiros para a geração de texto. Por omissão, e por razões de desempenho
(\Cref{sec:impl_m3}), o motor de texto na nuvem é o Groq e o Gemini serve a
visão e a rede de segurança final; mas o caminho totalmente local
permanece um \emph{fallback} de primeira classe, garantindo que o sistema
funciona --- e pode ser privado --- mesmo offline.

\section{\emph{Deployment}}
\label{sec:impl_deploy}

A solução é implantada em três serviços geridos: \textbf{Vercel}
(\emph{frontend} estático, com reescrita de rotas desconhecidas para
\texttt{index.html}); \textbf{Render} (\emph{backend} FastAPI a partir de
um contentor \textbf{Docker} próprio, que instala FFmpeg, libmagic e
Tesseract; configuração declarada em \texttt{render.yaml}, com sonda em
\texttt{/healthz}); e \textbf{Supabase} (autenticação, PostgreSQL e
Storage). No ambiente gratuito, o ChromaDB e o \emph{worker} Celery são
desativados — possível precisamente graças às decisões de degradação
graciosa descritas acima. A configuração CORS suporta origens exatas e uma
expressão regular para os \emph{deploys} de pré-visualização do Vercel.

% =====================================================================
%  CAPÍTULO 7 — TESTES E VALIDAÇÃO
% =====================================================================
\chapter{Testes e Validação}
\label{cap:testes}

A validação do sistema combina três níveis: testes automatizados do
\emph{backend}, verificação estática por tipagem, e validação funcional do
\emph{pipeline} completo. Discutem-se ainda as métricas relevantes para
uma avaliação académica da qualidade dos resultados.

\section{Plano de Testes}
\label{sec:plano_testes}

A estratégia de testes assentou em três pilares complementares:

\begin{enumerate}[leftmargin=*]
  \item \textbf{Testes unitários} — verificam, isoladamente, a lógica de
  componentes críticos e determinísticos.
  \item \textbf{Testes de integração} — verificam o comportamento de
  \emph{endpoints} da API, incluindo autenticação e sondas de saúde.
  \item \textbf{Verificação estática} — a tipagem (Python \emph{type
  hints} e TypeScript estrito) atua como uma camada de validação contínua,
  capturando incompatibilidades de esquema entre \emph{frontend} e
  \emph{backend} em tempo de compilação.
\end{enumerate}

\section{Testes Automatizados (pytest)}
\label{sec:pytest}

A suíte de testes do \emph{backend} foi implementada com \textbf{pytest}
(com \texttt{pytest-asyncio} e \texttt{pytest-cov}) e conta atualmente com
\textbf{37 testes a passar}. A \Cref{tab:testes} resume a sua distribuição.

\begin{table}[H]
\centering
\caption{Distribuição da suíte de testes automatizados.}
\label{tab:testes}
\small
\begin{tabularx}{\linewidth}{@{}llp{6.6cm}@{}}
\toprule
\textbf{Tipo} & \textbf{Módulo de teste} & \textbf{Âmbito} \\
\midrule
Unitário    & \texttt{test\_security}        & \emph{Security scanner}: \emph{checksum} e deteção de ficheiros inseguros. \\
Unitário    & \texttt{test\_upload\_validator} & Validação de tipo, \emph{magic bytes} e limites de tamanho. \\
Unitário    & \texttt{test\_date\_resolver}  & Resolução de datas incompletas e ambíguas (M2). \\
Unitário    & \texttt{test\_family\_graph}   & Construção do grafo e resumos de parentesco (M2). \\
Unitário    & \texttt{test\_templates}       & Integridade e formatação dos \emph{templates} narrativos (M3). \\
Unitário    & \texttt{test\_video\_builder}  & Funções de enquadramento e montagem de vídeo (M4). \\
Integração  & \texttt{test\_auth\_routes}    & Fluxo de autenticação na API. \\
Integração  & \texttt{test\_health\_route}   & Sondas de \emph{liveness} e \emph{deep health}. \\
\bottomrule
\end{tabularx}
\end{table}

A \Cref{fig:pytest} apresenta a execução da suíte. Os testes incidem
deliberadamente sobre a lógica \emph{determinística} (resolução de datas,
grafo, validação, \emph{templates}): a saída dos \textsc{llm} é
não-determinística e, por isso, validada por outras vias
(\Cref{sec:qualidade}).

\begin{figure}[H]
\centering
% \includegraphics[width=0.9\linewidth]{pytest.png}
\figph[4.5cm]{Captura de ecrã da execução \texttt{pytest -v} mostrando
os 37 testes a passar (idealmente com o sumário de cobertura
\texttt{--cov}).}
\caption{Resultado da execução da suíte de testes automatizados.}
\label{fig:pytest}
\end{figure}

\section{Validação Funcional do \emph{Pipeline}}
\label{sec:validacao_funcional}

Para além dos testes automatizados, validou-se o \emph{pipeline} completo
ponta-a-ponta com material de teste — incluindo ficheiros GEDCOM de
exemplo de várias famílias. O percurso \emph{happy path} validado foi:

\begin{enumerate}[leftmargin=*]
  \item Registo e autenticação de um utilizador.
  \item Carregamento de um conjunto de fotografias (M1: \textsc{exif} +
  análise visual + \emph{Storage}).
  \item Importação de uma árvore GEDCOM (M1/M2: pessoas + grafo).
  \item Geração de uma narrativa com tema, tom e idioma escolhidos (M3).
  \item Geração do vídeo documental a partir da narrativa (M4).
  \item Reprodução e descarregamento do vídeo final.
\end{enumerate}

Em cada etapa verificaram-se o estado das tarefas, as notificações ao
utilizador e a coerência dos artefactos produzidos.

\section{Métricas de Qualidade Narrativa}
\label{sec:qualidade}

A avaliação da qualidade da narrativa — intrinsecamente subjetiva — é um
desafio reconhecido (\Cref{sub:desafios}). Definem-se duas vias
complementares, em parte já instrumentáveis com os dados que o sistema
persiste (\texttt{task\_records}, \texttt{stories.facts\_used}):

\begin{itemize}[leftmargin=*]
  \item \textbf{Métricas objetivas.}
  \begin{itemize}
    \item \emph{Aderência ao PT-PT}: contagem automática de
    ``brasileirismos'' proibidos (a lista existe em \texttt{PT\_PT\_RULES})
    no texto gerado, produzindo uma percentagem de violações.
    \item \emph{Desempenho}: tempo médio de geração de narrativa e de
    vídeo, e taxa de sucesso/falha por \emph{template} e por \emph{backend}
    de \textsc{llm} (Ollama vs.\ Gemini).
  \end{itemize}
  \item \textbf{Métricas subjetivas.} Avaliação por um painel humano
  (escala 1--5) das dimensões \emph{coerência}, \emph{fidelidade factual}
  e \emph{qualidade literária}, sobre uma amostra de narrativas geradas.
\end{itemize}

A \Cref{tab:metricas} propõe a grelha de avaliação. A recolha sistemática
destes dados e a sua apresentação num painel constituem trabalho em curso,
discutido no \Cref{cap:conclusoes}.

\begin{table}[H]
\centering
\caption{Grelha proposta para avaliação da qualidade das narrativas
geradas.}
\label{tab:metricas}
\small
\begin{tabularx}{\linewidth}{@{}p{3.6cm}Yc@{}}
\toprule
\textbf{Métrica} & \textbf{Como medir} & \textbf{Tipo} \\
\midrule
Aderência PT-PT       & \% de violações às regras \texttt{PT\_PT\_RULES} & Objetiva \\
Fidelidade factual    & \% de factos da narrativa presentes no material   & Mista \\
Coerência             & Pontuação humana 1--5                              & Subjetiva \\
Qualidade literária   & Pontuação humana 1--5                              & Subjetiva \\
Tempo de geração      & Segundos (narrativa / vídeo)                       & Objetiva \\
Taxa de sucesso       & \% de gerações concluídas por \emph{template}      & Objetiva \\
\bottomrule
\end{tabularx}
\end{table}

\section{Resultados e Bugs Resolvidos}
\label{sec:resultados_testes}

Durante o desenvolvimento foram identificados e corrigidos diversos
defeitos relevantes, dos quais se destacam: a desalinhamento de nomes de
campos entre \emph{frontend} e \emph{backend} (por exemplo,
\texttt{narrative} vs.\ \texttt{content}, e \texttt{completed} vs.\
\texttt{ready} no estado dos vídeos), captado e resolvido graças à tipagem
estrita; a questão das tarefas eternamente ``em execução'', resolvida com
o \emph{sweep} de órfãs (\Cref{lst:sweep}); e a migração de
\texttt{datetime.utcnow()} para \texttt{datetime.now(UTC)} em conformidade
com as versões recentes do Python. Globalmente, o sistema cumpre os
requisitos funcionais essenciais (\rf{1}--\rf{13}), produzindo
narrativas e vídeos coerentes a partir de arquivos familiares reais.

% =====================================================================
%  CAPÍTULO 8 — CONCLUSÕES E TRABALHO FUTURO
% =====================================================================
\chapter{Conclusões e Trabalho Futuro}
\label{cap:conclusoes}

\section{Síntese e Cumprimento dos Objetivos}
\label{sec:sintese}

Este projeto demonstrou a viabilidade de um sistema completo que
transforma um arquivo familiar heterogéneo e desorganizado em narrativas
escritas e vídeos documentais coerentes, através de um \emph{pipeline}
modular de quatro estágios assente em inteligência artificial generativa.

Revendo os objetivos definidos no \Cref{sec:objetivos}: a ingestão
multimodal (O1), a organização temporal (O2), a geração narrativa (O3) e
a produção de vídeo (O4) foram \textbf{implementadas e validadas}; a
disponibilização através de uma aplicação Web moderna e implantada na
nuvem (O5) foi \textbf{concretizada} (FastAPI + React + Supabase + Render +
Vercel); e a premissa de privacidade e baixo custo (O6) foi
\textbf{assegurada} pela estratégia de IA híbrida e pelo uso de
ferramentas de código aberto. Em suma, os objetivos centrais do projeto
foram alcançados.

Para além dos objetivos definidos, o trabalho concretizou vários avanços
que elevam o resultado: a \textbf{narrativa segmentada em cenas}, que
sincroniza cada fotografia com a sua narração no vídeo (transformando o
\emph{slideshow} num verdadeiro documentário); um \textbf{executor de
tarefas \emph{in-process}} que substitui um \emph{worker} Celery pago no
plano gratuito (usado, por exemplo, na análise diferida das fotografias),
com limpeza de tarefas órfãs no arranque; uma \textbf{árvore familiar
interativa} (React~Flow) com pesquisa, destaque de linhagem, exportação em
imagem e edição/adição de familiares; legendas de vídeo como \textbf{faixa
WebVTT} sincronizada palavra-a-palavra, com escolha de voz; e ferramentas
de \textbf{curadoria de histórias} (favoritos, pesquisa/filtro,
regeneração orientada por \emph{feedback} e exportação para PDF). A
interface está totalmente \textbf{bilingue} (português europeu e inglês),
reforçada por um pós-processador determinístico que assegura a variante
europeia mesmo quando o modelo recai no português do Brasil.

Mais do que a integração de modelos, o trabalho destacou-se pela
\textbf{engenharia de robustez}: a degradação graciosa do RAG, os
\emph{fallbacks} de IA, o \emph{sweep} de tarefas órfãs, a execução
assíncrona \emph{in-process} e a recuperação de \emph{cold starts}
permitem ao sistema operar de forma fiável quer numa máquina local com
todas as dependências, quer num ambiente de nuvem restrito. A migração de
um protótipo \emph{local-first} para uma arquitetura multi-inquilino,
preservando a opção de inferência local, ilustra um compromisso maduro
entre acessibilidade, privacidade e escalabilidade.

\section{Limitações}
\label{sec:limitacoes}

Em honestidade técnica, reconhecem-se as seguintes limitações da versão
atual:

\begin{itemize}[leftmargin=*]
  \item \textbf{Avaliação quantitativa incompleta.} Embora os dados de
  base sejam persistidos, falta ainda um painel consolidado de métricas
  (tempos de geração, taxa de sucesso por \emph{template}, qualidade do
  RAG) e um protocolo de avaliação humana sistemático.
  \item \textbf{RAG inativo em produção.} O passo de \emph{retrieval} já
  está ligado à geração (\Cref{sub:geracao_narrativa}), mas no \emph{deploy}
  gratuito a ausência de ChromaDB força o modo \emph{stub} — só em ambiente
  local o \emph{retrieval} fica plenamente ativo. A migração para
  \emph{pgvector} (\Cref{sec:trabalho_futuro}) resolveria isto.
  \item \textbf{Sincronização de vídeo aproximada.} A duração de cada
  fotografia segue a narração da sua cena, mas pequenas derivas podem
  acumular-se em histórias longas; alinhar a troca de fotografias ao nível
  da palavra --- aproveitando os carimbos temporais por palavra já
  capturados para as legendas --- seria mais preciso.
  \item \textbf{Sem identificação de pessoas.} O reconhecimento facial
  (que tornaria as narrativas menos genéricas) não foi implementado.
  \item \textbf{Cobertura de testes do \emph{frontend}.} A automação de
  testes incide sobre o \emph{backend}; faltam testes ponta-a-ponta
  (E2E) da interface.
  \item \textbf{Acessibilidade não auditada formalmente.} Existem boas
  práticas, mas falta uma auditoria \textsc{wcag}~AA.
  \item \textbf{Geração de vídeo limitada pela memória no \emph{deploy}
  gratuito.} A montagem com MoviePy/FFmpeg a 720p excede os 512~MB de RAM do
  plano gratuito do Render, que termina a instância (\emph{out-of-memory}) a
  meio do \emph{render}. A solução adotada é uma \textbf{política de geração
  exclusivamente local}: na nuvem, o \emph{backend} \textbf{recusa} o
  \emph{render} (\texttt{VIDEO\_LOCAL\_ONLY}), evitando o
  \emph{out-of-memory}; os vídeos são pré-gerados em ambiente local (a 720p,
  sem limite de memória), escritos no \emph{mesmo} \emph{Storage} (Supabase)
  e reproduzidos no \emph{site} sem qualquer \emph{render} em tempo real ---
  a reprodução é leve (apenas servir o ficheiro) e nunca é afetada. A
  resolução e os \emph{fps} permanecem parametrizáveis por ambiente como
  ajuste adicional.
  \item \textbf{Quota de visão limitada no nível gratuito.} A análise visual
  das fotografias (M1) depende do Gemini, cujo nível gratuito permite apenas
  poucas dezenas de pedidos por dia; a geração de texto (M3) não é afetada por
  recorrer ao Groq como alternativa. A re-análise manual de fotografias é, por
  isso, parcimoniosa (só processa fotografias ainda sem descrição e confirma o
  custo em pedidos antes de avançar).
\end{itemize}

\section{Trabalho Futuro}
\label{sec:trabalho_futuro}

Identificam-se as seguintes direções de trabalho futuro, por ordem de
impacto esperado:

\begin{enumerate}[leftmargin=*]
  \item \textbf{Implantação local privada.} Executar o sistema completo numa
  máquina dedicada (com \textsc{gpu}) a correr um \textsc{llm} local via
  Ollama, eliminando a dependência de modelos na nuvem. \emph{Exemplo:} uma
  família gera as histórias no seu próprio servidor, sem que os dados saiam de
  casa.

  \item \textbf{Reconhecimento facial.} Identificar automaticamente as
  pessoas nas fotografias, com confirmação do utilizador, para tornar as
  narrativas mais específicas. \emph{Exemplo:} construir a linha temporal de
  uma pessoa a partir das fotografias em que é reconhecida.

  \item \textbf{Consulta conversacional do arquivo.} Permitir perguntas em
  linguagem natural sobre o arquivo, respondidas pelo subsistema \textsc{rag}.
  \emph{Exemplo:} responder a ``que pessoas aparecem nas fotografias de
  1985?''.

  \item \textbf{Mapa geo-temporal.} Posicionar as fotografias num mapa a
  partir das coordenadas \textsc{gps} do \textsc{exif}. \emph{Exemplo:}
  visualizar o percurso geográfico da família ao longo dos anos.

  \item \textbf{Ingestão de áudio e vídeo.} Estender o módulo de ingestão a
  vídeos e gravações de voz, transcritos automaticamente (e.g.\ com o
  \emph{Whisper}). \emph{Exemplo:} incorporar o relato gravado de um familiar
  na narrativa.

  \item \textbf{Clonagem de voz.} Narrar com a voz de um familiar a partir de
  uma amostra de áudio, mediante consentimento. \emph{Exemplo:} um
  documentário narrado na voz do próprio.

  \item \textbf{Avaliação formal.} Validar a qualidade das narrativas com um
  painel humano e métricas automáticas. \emph{Exemplo:} pontuar a coerência e
  a fidelidade factual numa amostra de histórias geradas.
\end{enumerate}

\section{Considerações Finais}
\label{sec:final}

O resultado deste projeto é uma plataforma funcional, segura e
estendível, que materializa uma ideia com valor humano evidente:
preservar e dar voz à memória familiar. Demonstra, em concreto, que é
hoje possível construir — com ferramentas de código aberto e a custo
essencialmente nulo — um sistema de IA generativa multimodal que
transforma um arquivo passivo num documentário narrado. O caminho de
evolução traçado, rumo a uma avaliação quantitativa rigorosa e a
funcionalidades de colaboração, confirma o potencial do trabalho para além
do âmbito desta licenciatura.


% ----------------------- ELEMENTOS PÓS-TEXTUAIS ---------------------
\cleardoublepage
\phantomsection
\addcontentsline{toc}{chapter}{Referências}
\bibliographystyle{ieeetr}
\bibliography{referencias}

\appendix
% =====================================================================
%  APÊNDICES
%  (Material próprio, demasiado extenso para o corpo do relatório.)
% =====================================================================
\chapter{Principais \emph{Endpoints} da API}
\label{ap:endpoints}

A \Cref{tab:endpoints} resume os principais \emph{endpoints} REST expostos
pelo \emph{backend} (prefixo \texttt{/api/v1/}).

\begin{longtable}{@{}p{1.4cm}p{6.6cm}p{5.4cm}@{}}
\caption{Principais \emph{endpoints} da API.}\label{tab:endpoints}\\
\toprule
\textbf{Método} & \textbf{Rota} & \textbf{Descrição} \\
\midrule
\endfirsthead
\toprule
\textbf{Método} & \textbf{Rota} & \textbf{Descrição} \\
\midrule
\endhead
GET    & \texttt{/healthz}                       & Sonda de saúde profunda. \\
POST   & \texttt{/upload} · \texttt{/upload/batch} & Carregamento de fotografias. \\
GET    & \texttt{/media}                         & Listar \emph{media}. \\
GET    & \texttt{/media/\{id\}/file}             & Servir ficheiro (auth por \emph{header} ou \texttt{?token=}). \\
POST   & \texttt{/genealogy/upload}              & Importar GEDCOM. \\
GET    & \texttt{/genealogy/persons}             & Listar pessoas. \\
GET    & \texttt{/timeline}                      & Linha temporal. \\
GET    & \texttt{/narrative/templates}           & Listar \emph{templates}. \\
POST   & \texttt{/narrative/generate}            & Gerar narrativa (\texttt{sync}/\texttt{background}). \\
GET    & \texttt{/narrative/stories}             & Listar histórias (favoritas primeiro). \\
PATCH  & \texttt{/narrative/stories/\{id\}}       & Editar história (título/texto/favorita). \\
POST   & \texttt{/narrative/stories/\{id\}/regenerate} & Regenerar a narrativa com \emph{feedback}. \\
POST   & \texttt{/multimedia/generate/\{id\}}    & Gerar vídeo a partir de uma história. \\
GET    & \texttt{/multimedia/videos}             & Listar vídeos. \\
GET    & \texttt{/multimedia/subtitle/\{ficheiro\}} & Servir a faixa de legendas \texttt{.vtt}. \\
GET/POST & \texttt{/projects} · \texttt{/projects/\{id\}} & CRUD de projetos. \\
GET    & \texttt{/tasks} · \texttt{/tasks/\{id\}} & Estado de tarefas assíncronas. \\
\bottomrule
\end{longtable}

\chapter{Configuração de \emph{Deployment}}
\label{ap:deploy}

Resumo da topologia de \emph{deployment} (ver \Cref{sec:impl_deploy}):

\begin{itemize}[leftmargin=*]
  \item \textbf{Vercel} — \emph{frontend} (SPA), \emph{build} Vite,
  reescrita SPA.
  \item \textbf{Render} — \emph{backend} via Docker (FFmpeg, libmagic,
  Tesseract), \texttt{render.yaml}, sonda \texttt{/healthz},
  \texttt{CELERY\_ENABLED=False} e ChromaDB desativado.
  \item \textbf{Supabase} — Auth (JWKS), PostgreSQL (asyncpg), Storage.
\end{itemize}

\noindent\emph{Nota:} as variáveis sensíveis (chaves Supabase, Groq e
Gemini) são definidas no painel do Render e nunca versionadas no
repositório.


\end{document}
